% !TeX program = lualatex
% !TeX encoding = utf8
% !TeX spellcheck = uk_UA
% !TeX root = ../main.tex

%=========================================================
\chapter{Електричний струм у~різних середовищах}\label{\currfilebase}\hypertarget{\currfilebase}{}
\graphicspath{{\currfilebase/Pictures}}

% --- Заглушки для рисунків: \figstub{ім'я-файлу.tikz}{що має бути зображено} ---
% Усі заглушки шукайте за \figstub; замінюйте на \localinput{ім'я-файлу.tikz}
\providecommand{\figstub}[2]{%
	\fbox{\parbox[c][4.2cm][c]{0.85\linewidth}{\centering\itshape\small
			[Заглушка рисунка \texttt{#1}]\\[2pt] #2}}}
%=========================================================

%=========================================================
\epigraph{Я не бачу способу уникнути висновку, що це~--- заряди негативної електрики, які переносяться частинками речовини. Постає
	питання: що це за частинки~--- атоми, молекули чи матерія в~ще тоншому стані поділу?
}{Дж.~Дж.~Томсон, \enquote{Катодні промені}, 1897}

Диференціальний закон Ома $\vect j = \lambda \Efield$~\eqref{eq:OhmDiff} є лише \emphx{загальним макроскопічним описом} провідності. Питома електропровідність $\lambda$ визначається мікроскопічною будовою середовища: концентрацією носіїв заряду, їхнім зарядом, рухливістю та механізмом взаємодії з~речовиною. Тому одна й~та сама напруженість електричного поля може спричиняти зовсім різні струми в~металі, електроліті, газі чи напівпровіднику.

Найважливіша відмінність між різними середовищами полягає в~\emphx{природі носіїв струму}. У~металах струм створюють вільні електрони; в~електролітах~--- позитивні та негативні йони; у~газах~--- електрони та йони, що виникають у~процесах іонізації; у~вакуумі струм створюється зарядженими частинками, емітованими електродами; у~напівпровідниках перенос заряду здійснюють електрони та дірки.

У~найзагальнішому випадку, якщо в~середовищі є кілька сортів носіїв, густина струму дорівнює сумі внесків усіх сортів (це узагальнення формули~\eqref{eq:JRhoU}; див. також відповідне зауваження в~гл.~\ref{SteadyCurrent}: внески носіїв протилежних знаків, що рухаються в~протилежних напрямках, \emphx{додаються}):
% ~~~~~~~~~~~~~~~~~~~~~~~~~~~~~~~~~~~~~~~~~~~~~~~~~~~~~~~~
\begin{equation}\label{eq:GeneralCurrentMedia}
\vect j = \sum_k \rho_k\vect u_k.
\end{equation}
% ~~~~~~~~~~~~~~~~~~~~~~~~~~~~~~~~~~~~~~~~~~~~~~~~~~~~~~~~
Тут $\rho_k=n_kq_k$~--- об'ємна густина заряду носіїв $k$-го сорту (не плутати з~питомим опором $\rho=1/\lambda$).

Якщо для кожного сорту носіїв виконується \emphx{лінійний закон}
% ~~~~~~~~~~~~~~~~~~~~~~~~~~~~~~~~~~~~~~~~~~~~~~~~~~~~~~~~
\begin{equation}\label{eq:LinearMobility}
	\vect u_k = \mu_k \vect E,
\end{equation}
% ~~~~~~~~~~~~~~~~~~~~~~~~~~~~~~~~~~~~~~~~~~~~~~~~~~~~~~~~
%
де $\mu_k$~--- \emphx{рухливість} носіїв, тобто дрейфова швидкість в~одиничному полі,
то, підставивши~\eqref{eq:LinearMobility} у~\eqref{eq:GeneralCurrentMedia}, дістаємо
% ~~~~~~~~~~~~~~~~~~~~~~~~~~~~~~~~~~~~~~~~~~~~~~~~~~~~~~~~
\begin{equation}\label{eq:GeneralConductivityMedia}
	\vect j = \sum_k n_k |q_k|\mu_k\,\vect E,
\end{equation}
% ~~~~~~~~~~~~~~~~~~~~~~~~~~~~~~~~~~~~~~~~~~~~~~~~~~~~~~~~
а~питома електропровідність
% ~~~~~~~~~~~~~~~~~~~~~~~~~~~~~~~~~~~~~~~~~~~~~~~~~~~~~~~~
\begin{equation}\label{eq:GeneralConductivity}
	\lambda = \sum_k n_k |q_k|\mu_k.
\end{equation}
% ~~~~~~~~~~~~~~~~~~~~~~~~~~~~~~~~~~~~~~~~~~~~~~~~~~~~~~~~

\begin{Attention}
	Співвідношення~\eqref{eq:LinearMobility}~--- не тотожність, а~\emphx{наближення слабкого поля}. Воно справедливе, доки дрейфова швидкість мала порівняно з~тепловою, $u_k\ll v_T$, і~час вільного пробігу $\tau$ можна вважати сталим. У~сильних полях $\tau$ починає залежати від енергії носія, і~залежність $u_k(E)$ стає нелінійною: у~напівпровідниках дрейфова швидкість насичується, у~газах виникає ударна іонізація, у~металах закон Ома порушується через нагрів. Саме тому закон Ома $\vect j=\lambda\vect E$~--- це \emphx{граничний} закон, а~не універсальний; у~цій главі ми побачимо середовища (газовий розряд, вакуумний діод, $p$-$n$-перехід), де він незастосовний.
\end{Attention}

Отже, макроскопічно різні провідники підкоряються одній і~тій самій схемі: \emphx{поле приводить у~впорядкований рух носіїв
заряду, а~властивості середовища визначають, наскільки легко цей рух виникає}.

% ========================================================
\begin{Attention}
\emphx{Зауваження про одиниці.} Більшість співвідношень цієї глави має однаковий вигляд у~гауссовій системі та в~СІ; там, де вигляд формули залежить від системи одиниць (закон $3/2$ для вакуумного діода, час релаксації заряду), наведено обидва варіанти. Числові значення (сталу Фарадея, число Лоренца, рухливості, питомі опори) подано в~СІ, як і~в~табл.~\ref{tab:ResistivityMetals}. Про розмірності електричних величин у~гауссовій системі~--- див. підрозділ <<Про одиниці вимірювання>> розд.~\hyperlink{ohms-law}{<<Закон Ома>>} гл.~\ref{SteadyCurrent}.
\end{Attention}
% ========================================================


%% --------------------------------------------------------
\section{Електричний струм у~металах}\label{sec:CurrentInMetals}
%% --------------------------------------------------------

У~металах частина валентних електронів слабко пов'язана з~окремими атомами й~утворює колективний електронний газ. Позитивні
йони утворюють кристалічну ґратку, відносно якої електрони можуть переміщуватися.

За відсутності зовнішнього поля електрони перебувають у~стані хаотичного руху. Середня швидкість цього руху дорівнює нулю,
тому макроскопічний струм відсутній:


% ========================================================
\begin{equation*}
		\langle\vect v\rangle = 0.
\end{equation*}
% ========================================================


Під дією електричного поля виникає мала добавка до хаотичного руху~--- дрейфова швидкість $\vect u$, і


% ========================================================
\begin{equation*}
		\vect j = -en\vect u_e,
\end{equation*}
% ========================================================
де $n$~--- концентрація вільних електронів, а~$\vect u_e$~--- їхня дрейфова швидкість.

Це окремий випадок формули~\eqref{eq:JRhoU} з~$\rho=-en$. Що саме електрони, а~не йони ґратки, є носіями струму в~металах, довели досліди Толмена і~Стюарта (розд.~\hyperlink{charge-carriers-nature}{<<Природа носіїв заряду>>} гл.~\ref{SteadyCurrent}); там же на прикладі срібного дроту показано, що дрейфова швидкість (частки міліметра за секунду) приблизно на дев'ять порядків менша за теплову ($\sim10^5$~м/с).

Оскільки електрон має від'ємний заряд, його дрейф відбувається проти напрямку електричного поля, але густина струму $\vect j$
спрямована за полем.

% \begin{Attention}
% У металі електрони не \enquote{рухаються від батарейки до лампи}. Їхній дрейф дуже повільний, тоді як електричне поле
% поширюється вздовж кола з електромагнітною швидкістю. Тому при замиканні кола впорядкований рух виникає практично
% одночасно по всій провідній системі.

% Саме поле в кожній ділянці дроту створюють, як пояснено в гл.~\ref{SteadyCurrent} (зауваження про <<естафету>> поляризації), \emphx{поверхневі заряди} цієї ділянки, а не безпосередньо джерело.
% \end{Attention}

\begin{Historical}
	Думка, що струм у~металі~--- це рух <<електронного газу>>, з'явилася майже одразу після відкриття електрона Дж.~Дж.~Томсоном (1897): у~1898~р. Е.~Рікке висунув гіпотезу вільних електронів у~металі, а~в~1900~р. П.~Друде побудував на її основі кількісну теорію. У~1905~р. Г.~Лоренц застосував до електронів газу статистику Максвелла--Больцмана й~кінетичне рівняння Больцмана (звідси назва <<модель Друде--Лоренца>>).
\end{Historical}

%% --------------------------------------------------------
\subsection*{Теорія провідності металів Друде}
%% --------------------------------------------------------

У~1900~р. П.~Друде запропонував просту класичну модель, що якісно (а~часто й~кількісно правильно за порядком величини) пояснює електропровідність
металів. У~цій моделі вільні електрони металу розглядають як класичний \enquote{газ}, що рухається в~полі $\Efield$ і~час від часу зазнає зіткнень з
іонами кристалічної ґратки, розсіюючись на них.

% \begin{figure}[h!]\centering
% 	\includegraphics[width=0.95\linewidth]{DrudeTrajectory}
% 	\caption{Хаотичний рух електрона в металі без поля та з полем: зіткнення з іонами ґратки й повільний дрейф\label{pic:DrudeTrajectory}}
% \end{figure}

Рівняння руху окремого електрона між зіткненнями (розглядаємо модулі величин, $e>0$~--- елементарний заряд): $a = eE/m_e$. Якщо одразу після зіткнення швидкість електрона (в~середньому)
дорівнює нулю, то до наступного зіткнення вона зростає лінійно з~часом, $v = at$ (рис.~\ref{pic:Ut}); після зіткнення набута впорядкована швидкість,
в~середньому, знову \enquote{обнуляється}, і~цикл прискорення повторюється.

\begin{figure}[h!]\centering
	\localinput{ut.tikz}
	\caption{Пилкоподібна залежність швидкості електрона від часу в~моделі Друде\label{pic:Ut}}
\end{figure}



Якщо середній час вільного пробігу дорівнює $\tau$, то середня (за час між зіткненнями) швидкість упорядкованого руху
% ~~~~~~~~~~~~~~~~~~~~~~~~~~~~~~~~~~~~~~~~~~~~~~~~~~~~~~~~
\begin{equation*}
	u = \overline v = \frac12\, \frac{e E \tau}{m_e},
\end{equation*}
% ~~~~~~~~~~~~~~~~~~~~~~~~~~~~~~~~~~~~~~~~~~~~~~~~~~~~~~~~
звідки, за формулою~\eqref{eq:JRhoU}, $j = enu = \dfrac{ne^2\tau}{2m_e} E$, тобто питома провідність за теорією Друде
% ~~~~~~~~~~~~~~~~~~~~~~~~~~~~~~~~~~~~~~~~~~~~~~~~~~~~~~~~
\begin{equation}\label{eq:DrudeConductivity}
	\tcbhighmath{
		\lambda = \frac{ne^2\tau}{2m_e}.
	}
\end{equation}
% ~~~~~~~~~~~~~~~~~~~~~~~~~~~~~~~~~~~~~~~~~~~~~~~~~~~~~~~~

\begin{Attention}
	Множник $\tfrac12$ у~формулі~\eqref{eq:DrudeConductivity}~--- наслідок \emphx{спрощеного} припущення, що всі електрони пролітають між
	зіткненнями рівно однаковий час $\tau$ (звідси й~\enquote{пилкоподібний} графік $v(t)$, лінійний на кожній ділянці, із середнім значенням рівно
	посередині). Насправді ж час вільного пробігу~--- випадкова величина з~(приблизно) експоненційним розподілом, і~акуратне усереднення за цим
	розподілом \emphx{прибирає} множник $1/2$, даючи $\lambda = ne^2\tau/m_e$, де $\tau$ тепер~--- \emphx{середній} час вільного пробігу.
	Розбіжність рівно вдвічі~--- класична \enquote{пастка} елементарного виведення теорії Друде, на яку варто звертати увагу: за порядком величини обидві
	формули однаково добре узгоджуються з~експериментом (сама теорія Друде~--- лише грубе класичне наближення, оскільки насправді електрони в
	металі підпорядковуються квантовій статистиці Фермі--Дірака), тому конкретний числовий множник тут не варто сприймати надто серйозно.
\end{Attention}

%% --------------------------------------------------------
\subsection*{Модель Друде--Лоренца}
%% --------------------------------------------------------

Як зазначено вище, множник $\tfrac12$ у~формулі Друде~\eqref{eq:DrudeConductivity} зникає, якщо коректно усереднити за розподілом
часів вільного пробігу. Уточнена питома провідність (теорія \emphx{Друде--Лоренца}):
% ~~~~~~~~~~~~~~~~~~~~~~~~~~~~~~~~~~~~~~~~~~~~~~~~~~~~~~~~
\begin{equation}\label{eq:DrudeLorentzConductivity}
	\tcbhighmath{
		\lambda_{\text{ДЛ}} = \frac{ne^2\tau}{m_e}.
	}
\end{equation}
% ~~~~~~~~~~~~~~~~~~~~~~~~~~~~~~~~~~~~~~~~~~~~~~~~~~~~~~~~
% Саме цю формулу зазвичай і мають на увазі під <<класичною теорією провідності металів>>.

Звідси видно, що провідність визначається двома мікроскопічними факторами: $\lambda\sim n\tau$. Велика концентрація вільних
електронів забезпечує металам високу провідність, а~зіткнення електронів з~ґраткою та її дефектами обмежують її.

Рухливість електронів у~металі, за означенням $\vect u=\mu\Efield$, у~моделі Друде--Лоренца дорівнює $\mu=e\tau/m_e$, так що $\lambda=ne\mu$ узгоджується із загальною формулою~\eqref{eq:GeneralConductivity}. Мікроскопічне виведення закону Джоуля--Ленца в~гл.~\ref{SteadyCurrent} (розд.~\hyperlink{joule-lenz-law}{<<Закон Джоуля--Ленца>>}) виконано з~множником $\tfrac12$; з~уточненим усередненням результат $w=\lambda E^2$ зберігається, тільки $\lambda$ набуває вигляду~\eqref{eq:DrudeLorentzConductivity}.

Час вільного пробігу $\tau$, а~отже, і~провідність $\lambda$, залежать від температури, оскільки з~підвищенням температури зростає інтенсивність
теплових коливань іонів ґратки, на яких розсіюються електрони. Дійсно, час вільного пробігу
$\tau = \ell/\bar v$, де $\ell$~--- довжина вільного пробігу (відстань між сусідніми зіткненнями), а~$\bar v$~--- швидкість електрона \emphx{між}
зіткненнями; оскільки дрейфова швидкість мізерна порівняно з~тепловою (гл.~\ref{SteadyCurrent}, розд.~\hyperlink{charge-carriers-nature}{<<Природа носіїв заряду>>}), $\bar v\approx v_T$. Найпростіший варіант теорії Друде
припускає, що концентрація носіїв $n$ і~довжина вільного пробігу $\ell$ від температури не залежать, а~сам електронний газ підпорядковується
розподілу Максвелла, тобто $\tfrac12 m_e v_T^2 \sim \tfrac32 k_BT$, звідки $v_T\propto\sqrt T$. Тоді
% ~~~~~~~~~~~~~~~~~~~~~~~~~~~~~~~~~~~~~~~~~~~~~~~~~~~~~~~~
\begin{equation*}
	\tau = \frac{\ell}{v_T} \propto \frac1{\sqrt T} \quad\Rightarrow\quad
	\lambda = \frac{ne^2\tau}{m_e} \propto \frac1{\sqrt T},
\end{equation*}
% ~~~~~~~~~~~~~~~~~~~~~~~~~~~~~~~~~~~~~~~~~~~~~~~~~~~~~~~~
а~отже, згідно з~\eqref{eq:DrudeLorentzConductivity}, для опору~--- оберненої до провідності величини:
% ~~~~~~~~~~~~~~~~~~~~~~~~~~~~~~~~~~~~~~~~~~~~~~~~~~~~~~~~
\begin{equation*}
	\rho = \frac1\lambda \propto \sqrt T
\end{equation*}
% ~~~~~~~~~~~~~~~~~~~~~~~~~~~~~~~~~~~~~~~~~~~~~~~~~~~~~~~~
(рис.~\ref{pic:DrudeRho}). Реальна залежність питомого опору чистих металів від температури для не надто низьких температур~--- майже
\emphx{лінійна}:
% ~~~~~~~~~~~~~~~~~~~~~~~~~~~~~~~~~~~~~~~~~~~~~~~~~~~~~~~~
\begin{equation}\label{eq:TempDependence1}
	\rho = \rho_0\left(1 + \alpha t\right)
\end{equation}
% ~~~~~~~~~~~~~~~~~~~~~~~~~~~~~~~~~~~~~~~~~~~~~~~~~~~~~~~~
де $t$~--- температура за шкалою Цельсія, а~$\rho_0$~--- питомий опір при $0\,^\circ$C, або, якщо відлічувати температуру від зручного початку відліку $20\,^\circ$C,
% ~~~~~~~~~~~~~~~~~~~~~~~~~~~~~~~~~~~~~~~~~~~~~~~~~~~~~~~~
\begin{equation}\label{eq:TempDependence2}
	\rho = \rho_{20}\left[1 + \alpha(t - 20)\right].
\end{equation}
% ~~~~~~~~~~~~~~~~~~~~~~~~~~~~~~~~~~~~~~~~~~~~~~~~~~~~~~~~

\begin{SCfigure}[2][h!]
	\centering
	% \begin{subfigure}[t]{0.45\linewidth}
		\centering
		\localinput{DrudeRho.tikz}
		\caption{Передбачення найпростішої теорії Друде: $\rho\propto\sqrt T$}
		\label{pic:DrudeRho}
	% \end{subfigure}\hfill
	% \begin{subfigure}[t]{0.45\linewidth}
	% 	\centering
	% 	\localinput{ExpRho.tikz}
	% 	\caption{Типова експериментальна залежність: лінійна ділянка й стрибок у нуль при
	% 		надпровідному переході}
	% 	\label{pic:ExpRho}
	% \end{subfigure}
	% \caption{Порівняння теоретичної моделі Друде та типової експериментальної
	% 	залежності питомого опору металів від температури}
	% \label{pic:RhoComparison}
\end{SCfigure}


У~табл.~\ref{tab:ResistivityMetals} наведено типові значення $\rho_{20}$ та $\alpha$ для кількох
поширених металів.

Оскільки опір $R$ у~законі Ома~\eqref{eq:OhmIntegral} залежить від температури, а~провідник зі струмом нагрівається (закон Джоуля--Ленца~\eqref{eq:JouleLenzIntegral}), при великих струмах вольт-амперна характеристика металу стає нелінійною. Так, опір вольфрамової спіралі лампи розжарювання в~розжареному стані ($\sim3000$~К) у~10--15 разів більший, ніж у~холодному, тому при вмиканні лампи струм короткочасно у~багато разів перевищує усталений.

\begin{table}[h!]\small\centering
	\caption{Питомий опір деяких металів при $20\,^\circ$C}
    \label{tab:ResistivityMetals}
		\begin{tblr}{
		colspec={X[l, m]X[c,m]X[c,m]},
        row{odd} = {gray!10},
		row{1} = {c, m, fg=white, bg=themecolorlight, font=\bfseries},
		hlines,
		vlines
		}
		Метал              & $\rho,\ 10^{-8}\,\text{Ом}\cdot\text{м}$ & $\alpha,\ 10^{-3}\,{}^\circ\text{C}^{-1}$ \\
		Мідь (\ce{Cu})     & 1.68                                     & 4.3                                       \\
		Алюміній (\ce{Al}) & 2.82                                     & 4.2                                       \\
		Залізо (\ce{Fe})   & 9.71                                     & 6.0                                       \\
		Нікель (\ce{Ni})   & 6.84                                     & 6.5                                       \\
		Вольфрам (\ce{W})  & 5.65                                     & 5.0                                       \\
	\end{tblr}
\end{table}

\begin{Attention}
	Лінійне зростання опору з~температурою характерне саме для \emphx{металів}, у~яких концентрація вільних носіїв $n$ практично не залежить від
	температури, а~опір визначається тільки частотою зіткнень. У~\emphx{напівпровідниках} залежність протилежна (розд.~\ref{sec:CurrentInSemiconductors}). Ще одне яскраве явище~--- \emphx{надпровідність} (див. далі).
\end{Attention}

%% --------------------------------------------------------
\subsection*{Квантова природа провідності металів}
%% --------------------------------------------------------

Класична теорія, навіть в~уточненому варіанті Друде--Лоренца, не може пояснити три
дослідні факти: лінійне зростання опору з~температурою (замість $\rho\propto\sqrt T$),
залишковий опір $\rho_0$ і~великі довжини вільного пробігу в~чистих металах.
Правильна теорія~--- квантова; для нашого викладу достатньо двох її ідей.

\emphx{Перша ідея: принцип Паулі.} Два електрони не можуть перебувати в~одному
квантовому стані, тому електрони заповнюють рівні знизу вгору. Енергію найвищого
заповненого рівня називають \emphx{енергією Фермі} $\varepsilon_F$ (для міді
$\approx7$~еВ, що в~сотні разів більше за $k_BT$ при $300$~К). Змінити стан можуть
лише електрони поблизу $\varepsilon_F$, тож у~переносі заряду бере участь лише
невелика їх частка, але рухаються вони зі швидкістю Фермі
% ~~~~~~~~~~~~~~~~~~~~~~~~~~~~~~~~~~~~~~~~~~~~~~~~~~~~~~~~
\begin{equation}\label{eq:FermiVelocity}
	v_F=\sqrt{\frac{2\varepsilon_F}{m_e}}\approx1{,}6\cdot10^{6}\ \text{м/с}
\end{equation}
% ~~~~~~~~~~~~~~~~~~~~~~~~~~~~~~~~~~~~~~~~~~~~~~~~~~~~~~~~
яка не залежить від температури. Тому час пробігу $\tau=\ell/v_F$ визначається лише
довжиною пробігу $\ell$, а~оскільки розсіяння на теплових коливаннях дає
$\ell\propto1/T$, дістаємо $\rho\propto T$~--- саме лінійну залежність, яку
спостерігають на досліді.

\emphx{Друга ідея: хвильова природа електрона.} Електрон поширюється в~кристалі як
хвиля і~проходить крізь \emphx{ідеально} періодичну ґратку без розсіяння, тому опір
такого кристала при $T=0$ дорівнює нулю. Опір виникає лише там, де періодичність
порушена: на теплових зміщеннях іонів та на домішках і~дефектах. Оскільки ці два
механізми діють незалежно, їхні внески додаються (\emphx{правило Маттіссена}):
% ~~~~~~~~~~~~~~~~~~~~~~~~~~~~~~~~~~~~~~~~~~~~~~~~~~~~~~~~
\begin{equation}\label{eq:Matthiessen}
	\tcbhighmath{
		\rho(T) = \rho_0 + \rho_{\text{фон}}(T),
	}
\end{equation}
% ~~~~~~~~~~~~~~~~~~~~~~~~~~~~~~~~~~~~~~~~~~~~~~~~~~~~~~~~
де $\rho_0$~--- \emphx{залишковий опір} (домішки й~дефекти), а~$\rho_{\text{фон}}(T)$
--- внесок теплових коливань ґратки (фононів). Другий доданок прямує до нуля при
$T\to0$, а~при високих температурах зростає лінійно. Нижче \emphx{температури Дебая}
$\Theta_D$ (для міді $\approx340$~К) він спадає швидше (приблизно як $T^5$, закон
Блоха--Грюнайзена), тому опір чистого металу виходить на плато $\rho_0$
(рис.~\ref{pic:Matthiessen}).

%---------------------------------------------------------
\begin{figure}[h!]\centering
    \localinput{Matthiessen.tikz}
%	\figstub{Matthiessen.tikz}{Залежність $\rho(T)$ для двох зразків одного металу (чистого й з домішками): криві зсунуті вздовж осі $\rho$ на $\Delta\rho_0$; біля нуля --- плато $\rho_0$ та ділянка $\propto T^5$, при високих $T$ --- лінійне зростання}
	\caption{Правило Маттіссена: домішки додають до опору сталу $\rho_0$, не змінюючи температурної частини\label{pic:Matthiessen}}
\end{figure}
%---------------------------------------------------------

\begin{Historical}
	Класичну модель вдалося виправити після появи квантової механіки: В.~Паулі (1925)
	сформулював принцип заборони, А.~Зоммерфельд (1927--1928) застосував його до
	електронів у~металі, а~Ф.~Блох (1928) показав, що електронні хвилі вільно
	поширюються в~ідеально періодичній ґратці. Адитивність опорів від різних причин
	розсіяння емпірично встановив А.~Маттіссен ще в~1860-х роках.
\end{Historical}

%% --------------------------------------------------------
\subsection*{Закон Відемана--Франца}
%% --------------------------------------------------------

Електрони в~металі переносять не лише заряд, а~й~тепло, тому теплопровідність
$\kappa$ і~електропровідність $\lambda$ мають бути пов'язані. Справді, за теорією
Друде--Лоренца обидві величини виражаються через ті самі $n$, $\tau$ і~$v_T$:
$\lambda=ne^2\tau/m_e$, $\kappa=\tfrac13nc_Vv_T^2\tau$, де $c_V=\tfrac32k_B$. Час
$\tau$ скорочується, і~залишається
% ~~~~~~~~~~~~~~~~~~~~~~~~~~~~~~~~~~~~~~~~~~~~~~~~~~~~~~~~
\begin{equation}\label{eq:WiedemannFranz}
	\frac{\kappa}{\lambda} = L T.
\end{equation}
% ~~~~~~~~~~~~~~~~~~~~~~~~~~~~~~~~~~~~~~~~~~~~~~~~~~~~~~~~
Величину $L$ називають \emphx{числом Лоренца} (на честь Л.~Лоренца, який у~1872~р.
виявив, що $\kappa/\lambda\propto T$; не плутати з~Г.~А.~Лоренцем, автором моделі
Друде--Лоренца). Класична теорія дає
$L_{\text{кл}}=\tfrac32(k_B/e)^2\approx1{,}1\cdot10^{-8}$~Вт$\cdot$Ом/К$^2$~\eqref{eq:LorenzClassical},
а~квантова~--- $L=\tfrac{\pi^2}{3}(k_B/e)^2\approx2{,}44\cdot10^{-8}$~Вт$\cdot$Ом/К$^2$~\eqref{eq:LorenzQuantum},
що добре узгоджується з~експериментом ($2{,}2$--$2{,}6\cdot10^{-8}$).

% ~~~~~~~~~~~~~~~~~~~~~~~~~~~~~~~~~~~~~~~~~~~~~~~~~~~~~~~~
\begin{equation}\label{eq:LorenzClassical}
	L_{\text{кл}} = \frac32\left(\frac{k_B}{e}\right)^2 \approx 1{,}1\cdot10^{-8}\ \text{Вт}\cdot\text{Ом}/\text{К}^2,
\end{equation}
% ~~~~~~~~~~~~~~~~~~~~~~~~~~~~~~~~~~~~~~~~~~~~~~~~~~~~~~~~
% ~~~~~~~~~~~~~~~~~~~~~~~~~~~~~~~~~~~~~~~~~~~~~~~~~~~~~~~~
\begin{equation}\label{eq:LorenzQuantum}
	L = \frac{\pi^2}{3}\left(\frac{k_B}{e}\right)^2 \approx 2{,}44\cdot10^{-8}\ \text{Вт}\cdot\text{Ом}/\text{К}^2,
\end{equation}
% ~~~~~~~~~~~~~~~~~~~~~~~~~~~~~~~~~~~~~~~~~~~~~~~~~~~~~~~~

\begin{Attention}
	Успіх моделі Друде в~поясненні закону Відемана--Франца був \emphx{випадковим}. Класична теорія помиляється двічі: вона вважає, що тепло переносять усі електрони, і~тому завищує теплоємність електронного газу приблизно у~сто разів; водночас вона використовує теплову швидкість $v_T$ замість значно більшої $v_F$~\eqref{eq:FermiVelocity}, і~квадрат швидкості занижує приблизно у~стільки ж разів. У~виразі для $\kappa$ ці помилки майже взаємно компенсуються, а~залишок і~дає різницю між~\eqref{eq:LorenzClassical} та~\eqref{eq:LorenzQuantum} (приблизно вдвічі). Сам Друде в~оригінальній роботі зробив ще одну помилку вдвічі й~отримав $L=3(k_B/e)^2\approx2{,}2\cdot10^{-8}$~--- майже точне експериментальне значення. Збіг з~дослідом ще не доводить правильності теорії.
\end{Attention}

%% --------------------------------------------------------
\subsection*{Надпровідність}
%% --------------------------------------------------------

У~1911~р. Г.~Камерлінг-Оннес виявив, що опір ртуті, охолодженої нижче критичної температури $T_c\approx4{,}15$~К, стрибкоподібно (а~не поступово)
звертається в~точний нуль (рис.~\ref{pic:superconductorR})~--- явище назвали \emphx{надпровідністю}. Це не просто <<дуже мала>> провідність: струм, одного
разу збуджений у~надпровідному кільці, циркулює без джерела ЕРС практично необмежено довго.

%---------------------------------------------------------
\begin{figure}[h!]\centering
\localinput{superconductorR.tikz}
\caption{Характерна залежність питомого опору від температури металу, який переходить в~надпровідний: при високих $T$~--- майже лінійне зростання, $T<T_c$ опір стрибкоподібно падає до нуля}
\label{pic:superconductorR}
\end{figure}
%---------------------------------------------------------

\begin{Historical}
	Відкриттю передувало зрідження гелію: у~1908~р. Х.~Камерлінг-Оннес у~Лейдені вперше
	отримав рідкий гелій при $\approx4{,}2$~К, а~відкачуючи пари~--- й~$1{,}7$~К, найнижчу
	на той час температуру. Це стало можливим завдяки каскадному охолодженню (хлористий
	метил, етилен, кисень, водень) і~подоланню головної трудності: гелій має надзвичайно
	низьку температуру інверсії, тож для ефекту Джоуля--Томсона його треба було попередньо
	охолодити рідким воднем. Саме гелієвий кріостат уможливив досліди 1911~р.
\end{Historical}

Нульовий опір змінює й~енергетику постійного струму. У~нормальному провіднику струм підтримується лише джерелом сторонніх сил, яке компенсує тепловтрати $\mathcal P=I^2R$~\eqref{eq:JouleLenzIntegral} (гл.~\ref{SteadyCurrent}). У~надпровідному кільці $R=0$, тепло не виділяється ($w=\vect j\cdot\Efield=0$, пор.~\eqref{eq:JouleLenzDiff}), і~збуджений струм тече без джерела ЕРС. Диференціальний закон Ома $\vect j=\lambda\Efield$ тут непридатний: $\lambda\to\infty$, $\Efield\to0$, а~густина струму залишається скінченною.

\emphx{Мікроскопічний механізм} надпровідності пояснює теорія
Бардіна--Купера--Шріффера (БКШ, 1957). Електрони відштовхуються кулонівськими
силами, але в~кристалі виникає слабке \emphx{притягання}, яке передає ґратка:
електрон, пролітаючи повз іони, деформує ґратку, і~до цієї деформації притягується
інший електрон (рис.~\ref{pic:CooperPair}). За низьких температур це притягання
об'єднує електрони в~\emphx{куперівські пари}, які поводяться як бозони й~збираються
в~єдиний колективний квантовий стан. Щоб розсіяти електрон у~такому стані, пару
треба зруйнувати, а~для цього потрібна енергія не менша за певну \emphx{енергетичну
щілину}; за низьких температур ґратка такої енергії не має, і~опір дорівнює нулю.

\begin{SCfigure}[2][h!]\centering
	\localinput{CooperPair.tikz}
	% \figstub{CooperPair.tikz}{Схема фононного механізму притягання: електрон деформує ґратку (згущення позитивних іонів), інший електрон притягується до цієї області; пара електронів з протилежними імпульсами та спінами}
	\caption{Утворення куперівської пари через обмін фононами\label{pic:CooperPair}}
\end{SCfigure}

Надпровідними стають далеко не всі метали. Мідь, срібло й~золото~--- найкращі
провідники при кімнатній температурі~--- не переходять у~надпровідний стан навіть
при охолодженні до мілікельвінів. Це не парадокс: та сама взаємодія електронів з
ґраткою, що визначає опір при кімнатній температурі, і~зв'язує електрони в~пари.
У~міді вона слабка~--- електрони рідко розсіюються на фононах, тому мідь добре
проводить струм, але притягання надто слабке для утворення пар. Навпаки, свинець і
ртуть мають при кімнатній температурі питомий опір у~десятки разів більший за мідь:
сильна взаємодія з~ґраткою заважає струму, але легко об'єднує електрони в~пари.
Тому \emphx{хороші провідники, як правило, погані надпровідники й~навпаки}.

% Куперівська пара має заряд $2e$. Це підтверджує \emphx{квантування магнітного потоку}: потік через надпровідне кільце може набувати лише значень, кратних $\Phi_0=h/(2e)\approx2{,}07\cdot10^{-15}$~Вб (СІ). Якщо ж розділити два надпровідники тонким шаром діелектрика, пари можуть <<просочуватися>> крізь нього без напруги (\emphx{ефект Джозефсона}, 1962); на ньому працюють чутливі магнітометри (SQUID) і надпровідні кубіти.

% За реакцією на магнітне поле розрізняють два типи надпровідників (рис.~\ref{pic:SCPhaseDiagram}). У надпровідниках \emphx{I роду} (чисті метали, наприклад ртуть, свинець) поле повністю виштовхується до критичного значення $H_c$, а при його перевищенні матеріал стрибком переходить у нормальний стан. У надпровідниках \emphx{II роду} (сплави й сполуки --- \ce{NbTi}, \ce{Nb3Sn}, купрати) є два критичні поля, $H_{c1}<H_{c2}$: між ними поле частково проникає у вигляді квантованих вихорів (вихорів Абрикосова), кожен з яких несе потік $\Phi_0$, а речовина в цілому залишається надпровідною. Саме тому матеріали II роду витримують поля в десятки тесла й використовуються в магнітах.

% \begin{figure}[h!]
% 	\centering
% 	\begin{subfigure}[t]{0.45\linewidth}
% 		\centering
% 		\figstub{SCTypeI.tikz}{Фазова діаграма $H(T)$ надпровідника I роду: крива $H_c(T)$ (парабола), нижче --- надпровідний стан}
% 		\caption{Надпровідник I роду}
% 	\end{subfigure}\hfill
% 	\begin{subfigure}[t]{0.45\linewidth}
% 		\centering
% 		\figstub{SCTypeII.tikz}{Фазова діаграма $H(T)$ надпровідника II роду: криві $H_{c1}(T)$ та $H_{c2}(T)$, між ними --- змішаний стан (вихори)}
% 		\caption{Надпровідник II роду}
% 	\end{subfigure}
% 	\caption{Фазові діаграми надпровідників у координатах магнітне поле -- температура\label{pic:SCPhaseDiagram}}
% \end{figure}


У~1986~р. Й.~Беднорц і~К.~Мюллер відкрили \emphx{високотемпературну надпровідність} у~купратній кераміці з~$T_c$, що згодом перевищила температуру
кипіння рідкого азоту ($77$~К)~--- це відкрило шлях до значно дешевшого азотного, а~не гелієвого, охолодження. Механізм високотемпературної
надпровідності досі остаточно не з'ясований: стандартна фононна теорія БКШ, найімовірніше, його не пояснює.

\begin{table}[h!]\small\centering
	\caption{Критична температура переходу деяких надпровідників}
	\label{tab:CriticalTemperatures}
	\begin{tblr}{
		colspec={X[l,m]X[c,m]X[l,m]},
        row{odd} = {gray!10},
		row{1} = {c, m, fg=white, bg=themecolorlight, font=\bfseries},
		hlines,
		vlines
	}
		Матеріал                            & $T_c$, К & Тип                          \\
		Ртуть (\ce{Hg})                     & 4.15     & звичайний (БКШ)              \\
		Свинець (\ce{Pb})                   & 7.2      & звичайний (БКШ)              \\
		Ніобій (\ce{Nb})                    & 9.25     & звичайний (БКШ)              \\
		\ce{NbTi} (сплав)                   & 10       & звичайний, технічний         \\
		\ce{Nb3Sn}                          & 18.3     & звичайний, технічний         \\
		\ce{MgB2}                           & 39       & проміжний                    \\
		\ce{YBa2Cu3O7} (YBCO)               & 93       & високотемпературний (купрат) \\
		\ce{Bi2Sr2Ca2Cu3O10} (BSCCO)        & 110      & високотемпературний (купрат) \\
		\ce{HgBa2Ca2Cu3O8}                  & 133      & високотемпературний (купрат) \\
		\ce{H3S} (під тиском $\sim150$~ГПа) & 203      & гідрид під високим тиском    \\
	\end{tblr}
\end{table}

У~табл.~\ref{tab:CriticalTemperatures} зібрано критичні температури кількох характерних надпровідників~--- від класичних металів, як-от ртуть
(перший виявлений надпровідник, $T_c\approx4{,}15$~К) чи ніобій, через проміжний випадок \ce{MgB2} ($39$~К) до купратних керамік, у~яких $T_c$
сягає $133$~К, і~гідридів під високим тиском, де вона перевищує $200$~К. Зверніть увагу на два орієнтири: \emphx{температура кипіння рідкого
гелію} ($4{,}2$~К)~--- межа, нижче якої працюють усі звичайні (низькотемпературні) надпровідники, і~\emphx{температура кипіння рідкого азоту}
($77$~К)~--- межа, вище якої охолодження стає значно дешевшим. Саме тому купрати з~$T_c>77$~К (першим був \ce{YBa2Cu3O7}, скорочено YBCO, $T_c=93$~К) називають \enquote{практично зручними}:
їх можна охолоджувати азотом, а~не дорогим гелієм.

Надпровідники застосовують там, де потрібні сильні магнітні поля або відсутність теплових втрат: обмотки надпровідних магнітів у~томографах (МРТ)
і~прискорювачах заряджених частинок (наприклад, диполі з~\ce{NbTi} у~Великому адронному колайдері, які працюють при $1{,}9$~К у~полі $8{,}3$~Тл, та котушки з~\ce{Nb3Sn} для ITER і~модернізації HL-LHC); лінії електропередач без омічних втрат; а~також надпровідні кубіти~--- один
з~провідних напрямів побудови квантових комп'ютерів.

%% --------------------------------------------------------
\section{Електричний струм в~електролітах}\label{sec:CurrentInElectrolytes}
%% --------------------------------------------------------

\emphx{Електролітами} називають речовини, водні або інші розчини яких проводять електричний струм завдяки наявності
рухливих йонів. До електролітів належать розчини кислот, лугів і~солей, а~також розплави багатьох йонних сполук.

На відміну від металів, у~електроліті струм переносять частинки \emphx{двох знаків}:
\begin{itemize}
\item позитивні йони~--- \emphx{катіони}, які рухаються до катода;
\item негативні йони~--- \emphx{аніони}, які рухаються до анода.
\end{itemize}

Нехай концентрація катіонів дорівнює $n_+$, їхній заряд $q_+>0$, а~середня дрейфова швидкість~--- $\vect u_+$.
Для аніонів відповідні величини позначимо $n_-$, $q_-<0$ та $\vect u_-$. Тоді
% ~~~~~~~~~~~~~~~~~~~~~~~~~~~~~~~~~~~~~~~~~~~~~~~~~~~~~~~~
\begin{equation}\label{eq:ElectrolyteCurrent}
\vect j = n_+q_+\vect u_+ + n_-q_-\vect u_-.
\end{equation}
% ~~~~~~~~~~~~~~~~~~~~~~~~~~~~~~~~~~~~~~~~~~~~~~~~~~~~~~~~

Це формула~\eqref{eq:JRhoU} для двох сортів носіїв протилежних знаків. Оскільки катіони рухаються за напрямом поля, а~аніони~--- проти нього, внески обох сортів у~густину струму мають один і~той
самий напрямок.

Для слабких полів дрейфову швидкість йонів можна записати у~вигляді

% ========================================================
\begin{equation*}
	\vect u_\pm = \mu_\pm \vect E,
\end{equation*}
% ========================================================


де $\mu_\pm$~--- рухливості відповідних йонів. Тоді
% ~~~~~~~~~~~~~~~~~~~~~~~~~~~~~~~~~~~~~~~~~~~~~~~~~~~~~~~~
\begin{equation}\label{eq:ElectrolyteConductivity}
\lambda = n_+q_+\mu_+ + n_-|q_-|\mu_-.
\end{equation}
% ~~~~~~~~~~~~~~~~~~~~~~~~~~~~~~~~~~~~~~~~~~~~~~~~~~~~~~~~

Рухливість йонів визначається в'язким тертям розчинника: за формулою Стокса, сила опору руху йона радіуса $a$ зі швидкістю $u$ дорівнює $6\pi\eta a u$ ($\eta$~--- в'язкість), тому усталений дрейф настає при $|q|E=6\pi\eta a u$, звідки $\mu=|q|/(6\pi\eta a)$. Типові рухливості йонів у~воді~--- порядку $10^{-8}$--$10^{-7}$~м$^2$/(В$\cdot$с); аномально висока рухливість \ce{H+} ($3{,}6\cdot10^{-7}$) пояснюється естафетною передачею протона між молекулами води (механізм Гроттгуса). З~підвищенням температури в'язкість води спадає, тому провідність електролітів, на відміну від металів, \emphx{зростає} з~температурою (приблизно на $2\%$ на кельвін). Від концентрації розчину $\lambda$ залежить немонотонно: спочатку зростає (більше йонів), а~в~концентрованих розчинах спадає (міжйонна взаємодія, неповна дисоціація).

На відміну від металів, проходження струму через електроліт супроводжується \emphx{перенесенням речовини}. Йони, рухаючись
до електродів, можуть брати участь у~хімічних реакціях. Це явище називають \emphx{електролізом} (рис.~\ref{pic:ElectrolysisCell}).

Струм у~колі неперервний: $\divg\vect j=0$~\eqref{eq:ContinuitySteady}. Це означає, що через будь-який переріз кола за одиницю часу проходить однаковий заряд. Але \emphx{носії} цього заряду в~різних ділянках різні: у~металевому електроді~--- електрони, в~електроліті~--- йони. На межі електрод--електроліт одні носії перетворюються на інші в~електродних реакціях: на катоді електрон з~металу приєднується до катіона, на аноді аніон віддає електрон металу. Саме тому сила струму однакова в~усіх перерізах кола, хоча природа носіїв уздовж кола змінюється.

\begin{Historical}
	Закони електролізу встановив М.~Фарадей (1833--1834); тоді ж він ввів терміни
	\emphx{електрод, анод, катод, електроліт, йон}. С.~Арреніус (1884--1887) показав,
	що електроліт розпадається на йони вже при розчиненні (Нобелівська премія з~хімії
	1903~р.). Закони електролізу дозволили Дж.~Стоні (1874) оцінити елементарний заряд
	$e=F/N_A$; він же запропонував назву <<електрон>> (1891).
\end{Historical}

\begin{figure}[h!]\centering
    \localinput{ElectrolysisCell.tikz}
	% \figstub{ElectrolysisCell.tikz}{Електролітична ванна: джерело, анод ($+$), катод ($-$), катіони до катода, аніони до анода, шар металу на катоді; стрілки $\vect j$ та напрямів дрейфу обох сортів йонів}
	\caption{Електроліз: рух йонів у~розчині й~осадження речовини на катоді\label{pic:ElectrolysisCell}}
\end{figure}

%% --------------------------------------------------------
\subsection*{Закони Фарадея для електролізу}
%% --------------------------------------------------------

Нехай через електроліт за час $t$ пройшов заряд

% ========================================================
\begin{equation*}
	q=It.
\end{equation*}
% ========================================================


Експериментально М.~Фарадей встановив, що маса речовини, виділена на електроді, пропорційна заряду, який пройшов через
електроліт:
% ~~~~~~~~~~~~~~~~~~~~~~~~~~~~~~~~~~~~~~~~~~~~~~~~~~~~~~~~
\begin{equation}\label{eq:FaradayFirst}
m=kq=kIt,
\end{equation}
% ~~~~~~~~~~~~~~~~~~~~~~~~~~~~~~~~~~~~~~~~~~~~~~~~~~~~~~~~
де $k$~--- \emphx{електрохімічний еквівалент} речовини.

\emphx{Другий закон Фарадея}: електрохімічний еквівалент пропорційний хімічному еквіваленту речовини, $k=\dfrac1F\cdot\dfrac{\mu}{z}$, де $\mu/z$~--- маса одного моля еквівалентів, а~$F$~--- універсальна стала.

Для одновалентних йонів маса одного моля речовини $\mu$ відповідає заряду $F$, який називають \emphx{сталою Фарадея}:
% ~~~~~~~~~~~~~~~~~~~~~~~~~~~~~~~~~~~~~~~~~~~~~~~~~~~~~~~~
\begin{equation}
F=N_Ae\approx9.65\cdot10^4\ \text{Кл/моль}.
\end{equation}
% ~~~~~~~~~~~~~~~~~~~~~~~~~~~~~~~~~~~~~~~~~~~~~~~~~~~~~~~~
Для йона з~валентністю $z$ заряд одного моля дорівнює $zF$. Тому другий закон Фарадея можна записати:
% ~~~~~~~~~~~~~~~~~~~~~~~~~~~~~~~~~~~~~~~~~~~~~~~~~~~~~~~~
\begin{equation}\label{eq:FaradaySecond}
\frac{m}{q}=\frac{\mu}{zF},
\end{equation}
% ~~~~~~~~~~~~~~~~~~~~~~~~~~~~~~~~~~~~~~~~~~~~~~~~~~~~~~~~
а~загальна формула для маси виділеної речовини набуває вигляду
% ~~~~~~~~~~~~~~~~~~~~~~~~~~~~~~~~~~~~~~~~~~~~~~~~~~~~~~~~
\begin{equation}\label{eq:ElectrolysisMass}
\tcbhighmath{
m=\frac{\mu}{zF}It.
}
\end{equation}
% ~~~~~~~~~~~~~~~~~~~~~~~~~~~~~~~~~~~~~~~~~~~~~~~~~~~~~~~~

\begin{Attention}
За виміряним зарядом $q$ та масою виділеної речовини $m$ можна визначити хімічний еквівалент речовини і, зрештою, кількість елементарних зарядів, що припадає на один йон.
\end{Attention}

На електродах можуть відбуватися складні хімічні реакції. Тому електроліз не слід розуміти просто як механічне
\enquote{осадження йонів}: електродні реакції визначають, яка саме речовина фактично утворюється на електроді.

%% --------------------------------------------------------
\subsection*{Електролітична комірка як неоднорідна ділянка кола}

Розчин з~електродами~--- це ділянка кола, на якій, окрім омічного спаду напруги, діє ще й~ЕРС хімічного походження. Якщо через комірку пропускають струм $I$ від зовнішнього джерела, то на межах електрод--розчин виникає \emphx{ЕРС поляризації} $\emf_{\text{пол}}$ (зворотна ЕРС: продукти електролізу самі утворюють гальванічний елемент, ЕРС якого напрямлена проти струму). Застосуємо закон Ома для неоднорідної ділянки~\eqref{eq:OhmNonUniform} з~$\phi_1-\phi_2=U$ і~$\emf=-\emf_{\text{пол}}$:
% ~~~~~~~~~~~~~~~~~~~~~~~~~~~~~~~~~~~~~~~~~~~~~~~~~~~~~~~~
\begin{equation}\label{eq:ElectrolysisOhm}
	IR = U - \emf_{\text{пол}}, \qquad I = \frac{U-\emf_{\text{пол}}}{R},
\end{equation}
% ~~~~~~~~~~~~~~~~~~~~~~~~~~~~~~~~~~~~~~~~~~~~~~~~~~~~~~~~
де $R$~--- опір розчину. Струм з'являється лише при $U>\emf_{\text{пол}}$ (<<напруга розкладання>>), а~потужність джерела $UI$ ділиться на джоулеве тепло в~розчині $I^2R$ (закон Джоуля--Ленца $w=\lambda E^2$ справедливий і~для йонної провідності, див.~\eqref{eq:JouleLenzDiff}) та корисну хімічну роботу $\emf_{\text{пол}}I$.

Обернений процес~--- \emphx{гальванічний елемент} та акумулятор: хімічні реакції на електродах відіграють роль \emphx{сторонніх сил} (див. підрозділ <<Сторонні сили. Електрорушійна сила>> гл.~\ref{SteadyCurrent}), які розділяють заряди й~створюють ЕРС $\emf$. Заряджання акумулятора~--- це електроліз з~$\emf_{\text{пол}}=\emf$, розряджання~--- зворотний хід тієї самої реакції.

%% --------------------------------------------------------
\section{Електричний струм у~газах}\label{sec:CurrentInGases}
%% --------------------------------------------------------

У~звичайному стані гази є дуже слабкими провідниками: більшість їхніх молекул та атомів електрично нейтральні, а~концентрація
вільних заряджених частинок дуже мала. Однак під дією зовнішнього чинника частина атомів або молекул може іонізуватися:

% ========================================================
\begin{equation*}
	\mathrm{A}\longrightarrow \mathrm{A}^+ + e^-.
\end{equation*}
% ========================================================


Внаслідок цього в~газі з'являються електрони та позитивні йони. За певних умов виникають також негативні йони.

\emphx{Іонізацією} називають процес утворення позитивних і~негативних заряджених частинок із нейтральних атомів або молекул.
Іонізацію газу можуть спричиняти нагрівання, ультрафіолетове та рентгенівське випромінювання, швидкі заряджені частинки,
радіоактивне випромінювання та сильне електричне поле.

Одночасно з~іонізацією відбувається зворотний процес~--- \emphx{рекомбінація}: позитивні та негативні носії зустрічаються
й~утворюють нейтральні частинки.

Отже, провідність газу визначається конкуренцією двох процесів:

% ========================================================
\begin{equation*}
	\text{іонізація}\quad\rightleftarrows\quad\text{рекомбінація}.
\end{equation*}

Кількісно, якщо іонізатор створює в~одиниці об'єму за одиницю часу $N$ пар, а~швидкість рекомбінації пропорційна добутку концентрацій носіїв обох знаків, $n_+n_-=n^2$, то за відсутності поля, що виносить носії на електроди, баланс носіїв описує рівняння
% ~~~~~~~~~~~~~~~~~~~~~~~~~~~~~~~~~~~~~~~~~~~~~~~~~~~~~~~~
\begin{equation}\label{eq:IonizationBalance}
	\frac{dn}{dt} = N - \beta n^2,
\end{equation}
% ~~~~~~~~~~~~~~~~~~~~~~~~~~~~~~~~~~~~~~~~~~~~~~~~~~~~~~~~
де $\beta$~--- коефіцієнт рекомбінації ($\beta\approx1{,}6\cdot10^{-12}$~м$^3$/с для йонів у~повітрі). Рівноважна концентрація $n_0=\sqrt{N/\beta}$; після вимкнення іонізатора концентрація спадає за законом $n=n_0/(1+\beta n_0t)$~--- не експоненційно. Рівняння~\eqref{eq:IonizationBalance}~--- це рівняння неперервності~\eqref{eq:ContinuityDiff} для носіїв одного сорту з~джерелом $N-\beta n^2$; для сумарного заряду джерела зникають, бо іонізація й~рекомбінація завжди відбуваються парами протилежних зарядів, і~закон збереження заряду виконується.
% ========================================================


%% --------------------------------------------------------
\subsection*{Несамостійний газовий розряд}
%% --------------------------------------------------------

Якщо зовнішній іонізатор постійно створює заряджені частинки, а~між електродами прикладено електричне поле, електрони та йони
починають рухатися впорядковано. Виникає \emphx{несамостійний газовий розряд}.

Нехай іонізатор створює $N$ пар йонів в~одиниці об'єму за одиницю часу; без поля концентрація носіїв визначається балансом~\eqref{eq:IonizationBalance}. При збільшенні напруженості поля струм спочатку зростає, оскільки все
більша частина створених заряджених частинок досягає електродів.

При достатньо великому полі практично всі створені іонізатором носії доходять до електродів. Подальше збільшення поля вже
майже не змінює струм: настає \emphx{струм насичення}.

Густину струму насичення знаходимо з~умови, що в~стаціонарному режимі кожна створена пара має дійти до електродів: у~плоскому конденсаторі з~відстанню $d$ між електродами
% ~~~~~~~~~~~~~~~~~~~~~~~~~~~~~~~~~~~~~~~~~~~~~~~~~~~~~~~~
\begin{equation}\label{eq:SaturationCurrentDensity}
	j_{\text{нас}} = eNd, \qquad I_{\text{нас}} = eNSd = eNV,
\end{equation}
% ~~~~~~~~~~~~~~~~~~~~~~~~~~~~~~~~~~~~~~~~~~~~~~~~~~~~~~~~
тобто струм насичення визначається \emphx{лише потужністю іонізатора} й~не залежить від напруги. При малих напругах, навпаки, носії встигають рекомбінувати раніше, ніж дійдуть до електродів, і~виконується закон Ома~\eqref{eq:OhmDiff} з
$\lambda=en_0(\mu_++\mu_-)=e(\mu_++\mu_-)\sqrt{N/\beta}$, де $\mu_\pm$~--- рухливості йонів. Ділянка закону Ома закінчується приблизно там, де омічна густина струму $\lambda E$
зрівнюється з~$j_{\text{нас}}$ (рис.~\ref{pic:GasIV}).

\begin{figure}[h!]\centering
	\localinput{GasDischargeIV.tikz}
	% \figstub{GasDischargeIV.tikz}{Схематична вольт-амперна характеристика розряду в газі: омічна ділянка, насичення, лавинне зростання (ділянка Таунсенда), пробій, тліючий розряд, перехід до дуги}
	\caption{Схематична вольт-амперна характеристика розряду в~газі: омічна ділянка
(1), перехідна ділянка (2), струм насичення (3), лавинне зростання й~пробій (4);
точки $a$--$e$~--- характерні режими\label{pic:GasIV}}
\end{figure}

На рис.~\ref{pic:GasIV} ділянки \emphx{1} і~\emphx{3}~--- це розглянуті вище омічна ділянка та струм насичення. На ділянці \emphx{2} зростання струму сповільнюється, бо частина носіїв ще встигає рекомбінувати. Ділянка \emphx{4}~--- лавинне зростання, де поле вже достатнє для ударної іонізації, і~розряд переходить у~самостійний (пробій, точка $e$); про нього йдеться далі.

%% --------------------------------------------------------
\subsection*{Самостійний газовий розряд}
%% --------------------------------------------------------

За достатньо сильного електричного поля сам газ може підтримувати свою іонізацію без зовнішнього іонізатора. Такий розряд
називають \emphx{самостійним}.

Важливу роль у~його виникненні відіграє \emphx{ударна іонізація}. Вільний електрон, прискорюючись електричним полем між
зіткненнями, може набути енергії, достатньої для іонізації нейтральної молекули:

% ========================================================
\begin{equation*}
	e^-+\mathrm{A}\longrightarrow 2e^-+\mathrm{A}^+.
\end{equation*}
% ========================================================


Число електронів при цьому збільшується. Нові електрони, у~свою чергу, також прискорюються полем і~можуть викликати подальшу
іонізацію. Виникає \emphx{електронна лавина} (рис.~\ref{pic:Avalanche}).

\begin{figure}[h!]\centering
	\localinput{ElectronAvalanche.tikz}
	% \figstub{ElectronAvalanche.tikz}{Електронна лавина між катодом та анодом: початковий електрон, ударна іонізація, число електронів зростає як $e^{\alpha_Tx}$; позитивні йони дрейфують до катода й вибивають вторинні електрони}
	\caption{Електронна лавина та вторинна емісія з~катода\label{pic:Avalanche}}
\end{figure}

Кількісно розмноження електронів описує \emphx{перший коефіцієнт Таунсенда} $\alpha_T$~--- число актів ударної іонізації, які здійснює один електрон на одиниці довжини шляху вздовж поля: на відстані $x$ від катода $n_e(x)=n_{e0}\,e^{\alpha_Tx}$, і~до анода (відстань $d$) із $n_{e0}$ електронів, що вийшли з~катода, долітає $n_{e0}e^{\alpha_Td}$. Позитивні йони, потрапляючи на катод, вибивають з~нього електрони (вторинна емісія): на кожен йон припадає $\gamma$ електронів (\emphx{другий коефіцієнт Таунсенда}). Розряд стає самостійним, коли ці вторинні електрони вже відновлюють початкову кількість, тобто виконується \emphx{умова Таунсенда}
% ~~~~~~~~~~~~~~~~~~~~~~~~~~~~~~~~~~~~~~~~~~~~~~~~~~~~~~~~
\begin{equation}\label{eq:TownsendCondition}
	\gamma\left(e^{\alpha_Td}-1\right)=1.
\end{equation}
% ~~~~~~~~~~~~~~~~~~~~~~~~~~~~~~~~~~~~~~~~~~~~~~~~~~~~~~~~
\emphx{Умова Таунсенда} визначає напругу запалювання. Коефіцієнт $\alpha_T$ залежить від $E/p$: довжина
вільного пробігу $\ell\propto1/p$, а~ймовірність набрати між зіткненнями енергію,
достатню для іонізації, $\propto\exp(-Bp/E)$. Тому $\alpha_T=Ap\exp(-Bp/E)$, і~з
умови Таунсенда дістаємо \emphx{закон Пашена} (1889):
% ~~~~~~~~~~~~~~~~~~~~~~~~~~~~~~~~~~~~~~~~~~~~~~~~~~~~~~~~
\begin{equation}\label{eq:Paschen}
	U_{\text{з}}=\frac{B\,pd}{\ln\dfrac{A\,pd}{\ln(1+1/\gamma)}}.
\end{equation}
% ~~~~~~~~~~~~~~~~~~~~~~~~~~~~~~~~~~~~~~~~~~~~~~~~~~~~~~~~
Напруга запалювання залежить лише від добутку $pd$ (рис.~\ref{pic:Paschen}) і~має
мінімум: при малих $pd$ електрон долітає до анода, майже не зустрічаючи молекул; при
великих $pd$ він часто стикається, не встигаючи набрати енергію. Для повітря
$U_{\min}\approx300$~В, а~електричну міцність при нормальному тиску оцінюють
$E_{\text{пр}}\approx3\cdot10^{6}$~В/м.

\begin{figure}[h!]\centering
	\localinput{Paschen.tikz}
	% \figstub{Paschen.tikz}{Крива Пашена: напруга запалювання $U_{\text{з}}$ як функція $pd$ (мінімум, різке зростання з обох боків)}
	\caption{Закон Пашена\label{pic:Paschen}}
\end{figure}

Різні режими самостійного газового розряду відрізняються умовами горіння, механізмом
підтримання та характерним світінням (рис.~\ref{pic:GasDischargeTypes},
табл.~\ref{tab:GasDischargeTypes}).

\begin{table}[h!]\small\centering
	\caption{Основні види самостійного газового розряду}
	\label{tab:GasDischargeTypes}
	\begin{tblr}{
		colspec={X[0.9,l,m]X[1.1,l,m]X[1.6,l,m]X[1.4,l,m]},
		row{odd} = {gray!10},
		row{1} = {c, m, fg=white, bg=themecolorlight, font=\bfseries},
		hlines,
		vlines
	}
		Вид & Умови & Механізм і особливості & Застосування \\
		Тліючий & $p\sim10^{-1}$--$10$~Торр, $U\sim10^2$~В, $I\sim$~мА &
			майже вся напруга падає в прикатодній області, де електрони
			розмножуються; далі --- слабко іонізований <<позитивний стовп>> &
			неонові трубки, люмінесцентні лампи, газові лазери \\
		Іскровий & атмосферний тиск, $E\gtrsim E_{\text{пр}}$ &
			короткочасний пробій вузькими каналами (стримерами); природний
			приклад --- блискавка ($I$ до $\sim10^{5}$~А) &
			іскрові розрядники, запалювання \\
		Коронний & сильно неоднорідне поле біля вістер і тонких проводів &
			поле біля електрода перевищує $E_{\text{пр}}$, далі слабке ---
			світиться лише тонка <<корона>>; природний прояв --- <<вогні
			святого Ельма>> &
			електрофільтри, копіювальні апарати; шкода --- втрати на ЛЕП \\
		Дуговий & $I\gtrsim1$~А, $U\sim10$--$50$~В &
			розжарений катод (термо- й автоелектронна емісія,
			розд.~\ref{sec:CurrentInVacuum}), температура каналу
			$\sim5000$--$10000$~К; ВАХ спадна, потрібен баласт &
			зварювання, дугові печі, прожектори \\
	\end{tblr}
\end{table}

Йонізований квазінейтральний газ, у~якому концентрації електронів і~йонів приблизно рівні, називають \emphx{плазмою} (термін запропонував І.~Ленгмюр, 1928). Плазмою є не лише розряд у~лампі, а~й~Сонце, і~полярне сяйво, і~майже вся видима речовина у~Всесвіті.

\begin{Historical}
	Дугу вперше спостерігав В.~Петров (1802). Теорію лавинного пробою розвинув
	Дж.~Таунсенд на початку XX~ст. В~Україні промислове дугове зварювання пов'язане з
	Інститутом електрозварювання імені Є.~О.~Патона (Київ, 1934).
\end{Historical}

\begin{figure}[h!]\centering
    % \localinput{GasDischargeTypes.tikz}
	\includegraphics[width=0.8\linewidth]{Discharges}
	% \figstub{GasDischargeTypes.jpg}{Фотографії чотирьох типів розряду в одному рисунку: іскра (блискавка), корона (вістря), тліючий (неонова трубка), дуга (зварювання)}
	\caption{Види самостійного газового розряду\label{pic:GasDischargeTypes}}
\end{figure}

\begin{Attention}
Газовий розряд принципово відрізняється від струму в~металі. У~металі кількість вільних носіїв приблизно задана будовою
речовини, тоді як у~газі концентрація носіїв може радикально змінюватися під час самого проходження струму. Тому в~газі
електричне поле може не просто переміщувати носії, а~й~змінювати їх кількість.
\end{Attention}


%% --------------------------------------------------------
\section{Електричний струм у~вакуумі}\label{sec:CurrentInVacuum}
%% --------------------------------------------------------

У~ідеальному вакуумі немає речовини, а~отже, немає і~готових носіїв струму. Проте струм у~вакуумі можливий, якщо джерело
заряджених частинок розміщене безпосередньо в~області вакууму.

Найважливішим механізмом є \emphx{електронна емісія}~--- вихід електронів із поверхні твердого тіла.

Залежно від механізму розрізняють:
\begin{itemize}
\item \emphx{термоелектронну емісію}~--- вихід електронів при нагріванні;
\item \emphx{фотоелектронну емісію}~--- вихід електронів під дією світла;
\item \emphx{автоелектронну емісію}~--- вихід електронів під дією дуже сильного електричного поля;
\item \emphx{вторинну емісію}~--- вибивання електронів із поверхні при бомбардуванні її швидкими частинками.
\end{itemize}

Усі ці механізми об'єднує одне: електрон не може вільно покинути метал, бо на його поверхні діє потенціальний бар'єр (рис.~\ref{pic:WorkFunction}). Мінімальну енергію, яку треба передати електрону для виходу, називають \emphx{роботою виходу} $A_{\text{вих}}$. Механізми емісії відрізняються лише тим, звідки електрон її отримує. При нагріванні~--- з~теплового руху. Під дією світла~--- з~енергії фотона: $h\nu=A_{\text{вих}}+\tfrac12m_ev^2$ (формула Ейнштейна, 1905; ефект спостерігав Г.~Герц у~1887~р.). У~дуже сильному полі ($\sim10^9$~В/м) бар'єр стає настільки тонким, що електрони <<просочуються>> крізь нього (тунельний ефект), навіть без нагріву. При вторинній емісії енергію передає швидка частинка. Типові значення $A_{\text{вих}}$: вольфрам~--- $4{,}5$~еВ, цезій~--- $2{,}1$~еВ, оксидні катоди (\ce{BaO})~--- близько $1$--$1{,}5$~еВ; саме тому оксидні катоди працюють при значно нижчих температурах.

У~найпростішому наближенні густина термоелектронного струму описується законом Річардсона--Дешмана:
% ~~~~~~~~~~~~~~~~~~~~~~~~~~~~~~~~~~~~~~~~~~~~~~~~~~~~~~~~
\begin{equation}\label{eq:RichardsonDushman}
	j = A_R T^2
	\exp\left(-\frac{A_{\text{вих}}}{k_BT}\right),
\end{equation}
% ~~~~~~~~~~~~~~~~~~~~~~~~~~~~~~~~~~~~~~~~~~~~~~~~~~~~~~~~
де $A_R=4\pi m_eek_B^2/h^3\approx1{,}2\cdot10^{6}$~А/(м$^2\cdot$К$^2$)~--- стала Річардсона (СІ; для реальних металів вона того ж порядку). Експоненціальний множник показує, що струм дуже чутливий до температури: його визначає частка електронів, чия енергія перевищує $A_{\text{вих}}$ (задача~\ref{prob:RichardsonRatio}).

\begin{figure}[h!]\centering
	\localinput{WorkFunction.tikz}
	% \figstub{WorkFunction.tikz}{Потенціальна яма електрона в металі: рівень Фермі, рівень вакууму, робота виходу $A_{\text{вих}}$; заповнені стани електронів}
	\caption{Потенціальна яма електрона в~металі: рівень Фермі, рівень вакууму, робота виходу $A_{\text{вих}}$; заповнені стани електронів\label{pic:WorkFunction}}
\end{figure}

\begin{Attention}
У~вакуумі струм не є \enquote{струмом вакууму} як властивістю самого простору. Вакуум лише не перешкоджає руху заряджених частинок. Провідність вакуумного проміжку визначається тим, скільки носіїв у~нього введено та як вони рухаються в~електричному полі.
\end{Attention}

Класичним прикладом вакуумного струму є електронна лампа (рис.~\ref{pic:VacuumDiode}). У~найпростішому вакуумному діоді нагрітий катод емітує електрони, а~позитивно заряджений анод притягує їх. За відповідної полярності анdода струм проходить, тоді як при зміні полярності електрони не можуть потрапити на анdод. Таким чином, вакуумний діod має властивість односторонньої провідності.

%---------------------------------------------------------
\begin{SCfigure}[1][h!]\centering
\includegraphics[width=0.45\linewidth]{Diod}
\caption{Вигяд вакуумного діода: нагрітий катод емітує електрони, які прискорюються до анода під дією електричного поля. При зміні полярності анода струм не проходить.}
\label{pic:VacuumDiode}
\end{SCfigure}
%---------------------------------------------------------

%% --------------------------------------------------------
\subsection*{Струм, обмежений просторовим зарядом}
%% --------------------------------------------------------

У~вакуумі носії рухаються без зіткнень, тому густина струму задається формулою~\eqref{eq:JRhoU}: $j=\rho u$, де $\rho$~--- густина заряду електронів, а~$u$~--- їхня швидкість (тепер не повільний дрейф у~речовині, а~швидкість вільного прискорення полем). У~стаціонарному режимі $\divg\vect j=0$~\eqref{eq:ContinuitySteady}, тож у~плоскому діоді $j=\rho(x)v(x)=\const$. 

Якщо катод емітує багато електронів, біля нього утворюється хмара просторового заряду, яка повністю екранує поле катода ($E(0) = -\left.\frac{\mathrm{d}\phi}{\mathrm{d}x}\right|_{x=0} = 0$) й~обмежує струм. Швидкість електрона на~відстані $x$ від катода визначається потенціалом $\phi(x)$ у~цій точці: $v=\sqrt{2e\phi/m_e}$ (із~закону збереження енергії), тому модуль густини заряду $|\rho| = j/v = j\sqrt{m_e/(2e)}\,\phi^{-1/2}$ (рис.~\ref{tikz:DiodeProblem}).

%---------------------------------------------------------
\begin{figure}[h!]\centering
\localinput{DiodeProblem.tikz}
\caption{Розподіл просторового заряду та електронів у плоскому вакуумному діоді.}\label{tikz:DiodeProblem}
\end{figure}
%---------------------------------------------------------

Підставимо $|\rho|$ у~одновимірне рівняння Пуассона $\frac{\mathrm{d}^2\phi}{\mathrm{d}x^2} = 4\pi|\rho|$:
\begin{equation*}
    \frac{\mathrm{d}^2\phi}{\mathrm{d}x^2} = 4\pi j \sqrt{\frac{m_e}{2e}} \, \phi^{-1/2}.
\end{equation*}
Помноживши обидві частини на~$2\frac{\mathrm{d}\phi}{\mathrm{d}x}$ та~проінтегрувавши з~урахуванням умови екранування на~катоді ($\phi(0)=0$, $\left.\frac{\mathrm{d}\phi}{\mathrm{d}x}\right|_{x=0} = 0$), маємо:
\begin{equation*}
    \left(\frac{\mathrm{d}\phi}{\mathrm{d}x}\right)^2 = 16\pi j \sqrt{\frac{m_e}{2e}} \, \phi^{1/2} 
    \implies 
    \frac{\mathrm{d}\phi}{\mathrm{d}x} = 4 \sqrt{\pi j} \left(\frac{m_e}{2e}\right)^{1/4} \phi^{1/4}.
\end{equation*}
Розділивши змінні й~інтегруючи від катода ($x=0, \phi=0$):
\begin{equation*}
    \int_0^{\phi(x)} \phi^{-1/4} \mathrm{d}\phi = 4 \sqrt{\pi j} \left(\frac{m_e}{2e}\right)^{1/4} \int_0^x \mathrm{d}x 
    \implies 
    \frac{4}{3} U^{3/4} = 4 \sqrt{\pi j} \left(\frac{m_e}{2e}\right)^{1/4} d.
\end{equation*}
Скоротивши обидві частини на~$4$, отримаємо зв'язок між потенціалом $\phi(x)$ у~довільній точці $x$ та~сталою густиною струму $j$:
\begin{equation*}
    \frac{1}{3} \phi^{3/4}(x) = \sqrt{\pi j} \left(\frac{m_e}{2e}\right)^{1/4} x.
\end{equation*}

Звідси потенціал у~проміжку зростає за~законом $\phi(x) \propto x^{4/3}$:
\begin{equation*}
    \phi(x) = \left( 3 \sqrt{\pi j} \right)^{4/3} \left(\frac{m_e}{2e}\right)^{1/3} x^{4/3}.
\end{equation*}

Застосуємо цей вираз до~всього проміжку (на~аноді $x = d$, $\phi(d) = U$):
\begin{equation*}
    \frac{1}{3} U^{3/4} = \sqrt{\pi j} \left(\frac{m_e}{2e}\right)^{1/4} d.
\end{equation*}

Піднісши цю рівність до~квадрата, маємо:
\begin{equation*}
    \frac{1}{9} U^{3/2} = \pi j \sqrt{\frac{m_e}{2e}}\, d^2.
\end{equation*}

Остаточно виражаємо шукану густину струму, що й~дає \emphx{закон трьох других} (закон Чайльда--Ленгмюра; у~вітчизняній літературі також~--- Богуславського--Ленгмюра):
% ~~~~~~~~~~~~~~~~~~~~~~~~~~~~~~~~~~~~~~~~~~~~~~~~~~~~~~~~
\begin{equation}\label{eq:ChildLangmuir}
    j = \frac{\sqrt2}{9\pi}\sqrt{\frac{e}{m_e}}\,\frac{U^{3/2}}{d^2}.
\end{equation}
% ~~~~~~~~~~~~~~~~~~~~~~~~~~~~~~~~~~~~~~~~~~~~~~~~~~~~~~~~

Струм росте не лінійно з~напругою, а~як $U^{3/2}$, тому закон Ома для вакуумного діода не виконується. Лише коли напруга така, що всі емітовані електрони відразу йдуть на анод, струм виходить на насичення, яке за законом Річардсона--Дешмана~\eqref{eq:RichardsonDushman} визначається температурою катода (рис.~\ref{pic:VacuumDiodeIV}). Електричний струм гріє не~простір між електродами, а~сам анод, у~який на~повній швидкості врізаються розігнані електрони, виділяючи потужність $\mathcal{P} = IU$.

\begin{figure}[h!]\centering
	\localinput{VacuumDiodeIV.tikz}
	% \includegraphics[width=0.35\linewidth]{Diod}
	% \figstub{VacuumDiodeIV.tikz}{Анодна характеристика вакуумного діода $I(U)$ для двох температур катода: ділянка $\propto U^{3/2}$ та плато насичення (вище для більшої $T$); праворуч --- схема діода}
	\caption{Вольт-амперна характеристика вакуумного діода\label{pic:VacuumDiodeIV}}
\end{figure}

\begin{Historical}
	Термоелектронну емісію вперше помітив Т.~Едісон (1880-ті, \emphx{ефект Едісона}).
	Дж.~Флемінг (1904) створив на цій основі вакуумний діод, а~Лі де Форест (1906)
	додав керівну сітку й~отримав тріод. Теорію термоемісії розвинув О.~Річардсон
	(Нобелівська премія 1928~р.).
\end{Historical}

%% --------------------------------------------------------
\section{Електричний струм у~напівпровідниках}\label{sec:CurrentInSemiconductors}
%% --------------------------------------------------------

\emphx{Напівпровідники} займають проміжне положення між металами та діелектриками за величиною електропровідності. До них
належать, зокрема, кремній, германій та багато складних напівпровідникових сполук.

Головна особливість напівпровідників полягає в~тому, що концентрація носіїв струму дуже чутлива до температури, домішок,
освітлення та інших зовнішніх чинників.

Причину такої поведінки пояснює \emphx{зонна теорія}. Енергії, які може мати електрон у~кристалі, утворюють смуги~--- \emphx{зони}, розділені забороненими інтервалами. Струм можливий лише тоді, коли електрон може змінити стан, тобто поруч є вільні стани (як і~для металів, див. принцип Паулі в~розд.~\ref{sec:CurrentInMetals}). У~металі верхня зона заповнена частково, тож вільні стани є завжди. У~діелектрику й~напівпровіднику при $T=0$ нижня (\emphx{валентна}) зона заповнена повністю, а~вища (\emphx{зона провідності}) порожня, тому струму немає. Різниця між ними лише в~ширині забороненої зони $E_g$. У~напівпровідників вона порівняно мала (кремній~--- $1{,}12$~еВ, германій~--- $0{,}66$~еВ, \ce{GaAs}~--- $1{,}42$~еВ), тому вже тепловий рух ($k_BT\approx0{,}026$~еВ при $300$~К) закидає в~зону провідності помітну кількість електронів, і~ця кількість швидко росте з~температурою. У~діелектрика ($E_g\gtrsim5$~еВ, алмаз~--- $5{,}5$~еВ) цього не відбувається (рис.~\ref{pic:Bands}, табл.~\ref{tab:Semiconductors}).

\begin{figure}[h!]\centering
	\localinput{Bands.tikz}
	% \figstub{Bands.tikz}{Три схеми енергетичних зон: метал (частково заповнена зона), напівпровідник (мала $E_g$) та діелектрик (велика $E_g$); виділено зону провідності, валентну зону, рівень Фермі}
	\caption{Зонна структура металу, напівпровідника й~діелектрика\label{pic:Bands}}
\end{figure}

\begin{table}[h!]\small\centering
	\caption{Ширина забороненої зони та рухливості носіїв у~деяких напівпровідниках при $300$~К}
	\label{tab:Semiconductors}
	\begin{tblr}{
		colspec={X[l,m]X[c,m]X[c,m]X[c,m]},
		row{odd} = {gray!10},
		row{1} = {c, m, fg=white, bg=themecolorlight, font=\bfseries},
		hlines,
		vlines
	}
		Матеріал                  & $E_g$, еВ & $\mu_n$, м$^2$/(В$\cdot$с) & $\mu_p$, м$^2$/(В$\cdot$с) \\
		Кремній (\ce{Si})         & 1.12      & 0.135                      & 0.048                      \\
		Германій (\ce{Ge})        & 0.66      & 0.39                       & 0.19                       \\
		Арсенід галію (\ce{GaAs}) & 1.42      & 0.85                       & 0.04                       \\
	\end{tblr}
\end{table}

У~власному напівпровіднику носіями струму є:
\begin{itemize}
\item \emphx{електрони}~--- частинки з~зарядом $-e$, що перебувають у~зоні провідності;
\item \emphx{дірки}~--- ефективні позитивні носії заряду, пов'язані з~відсутністю електрона у~валентній зоні.
\end{itemize}

Дірка не є самостійною елементарною частинкою. Це зручний спосіб опису колективного руху електронів у~майже заповненій
валентній зоні. Якщо сусідній електрон переходить на місце вакансії, сама вакансія переміщується в~протилежному напрямку,
ніби позитивно заряджена частинка (рис.~\ref{pic:HoleMotion}).

Електрон у~кристалі рухається в~періодичному полі ґратки, тому на зовнішнє поле він реагує так, ніби має іншу масу~--- \emphx{ефективну} $m^*$, яка може відрізнятися від маси вільного електрона. Рухливість при цьому $\mu=e\tau/m^*$, як і~в~металах (розд.~\ref{sec:CurrentInMetals}). Наприклад, у~\ce{GaAs} електрони легкі ($m^*\approx0{,}07\,m_e$), тому їхня рухливість висока (табл.~\ref{tab:Semiconductors}).

\begin{figure}[h!]\centering
	\localinput{HoleMotion.tikz}
	% \includegraphics[width=0.5\linewidth]{HoleMotion}
	% \figstub{HoleMotion.tikz}{Послідовність з трьох кадрів: ряд атомів з електронами, вакансія (дірка) зміщується проти руху електронів, тобто за напрямом поля}
	\caption{Рух дірки як естафета електронів на сусідні вакансії. Електрони позначені чорними кружками, дірки~--- білими. Дірки рухаються праворуч, електрони~--- ліворуч.\label{pic:HoleMotion}}
\end{figure}

Тому густина струму в~напівпровіднику має два внески:
% ~~~~~~~~~~~~~~~~~~~~~~~~~~~~~~~~~~~~~~~~~~~~~~~~~~~~~~~~
\begin{equation}\label{eq:SemiconductorCurrent}
\vect j =
-e n\,\vect u_n
+e p\,\vect u_p,
\end{equation}
% ~~~~~~~~~~~~~~~~~~~~~~~~~~~~~~~~~~~~~~~~~~~~~~~~~~~~~~~~
де $n$ і~$p$~--- концентрації електронів та дірок, а~$\vect u_n$ і~$\vect u_p$~--- їхні дрейфові швидкості.

У~лінійному режимі

% ========================================================
\begin{equation*}
		\vect u_n=-\mu_n\vect E,
		\qquad
		\vect u_p=\mu_p\vect E,
\end{equation*}
% ========================================================


тому
% ~~~~~~~~~~~~~~~~~~~~~~~~~~~~~~~~~~~~~~~~~~~~~~~~~~~~~~~~
\begin{equation}\label{eq:SemiconductorConductivity}
\tcbhighmath{
\lambda=e(n\mu_n+p\mu_p).
}
\end{equation}
% ~~~~~~~~~~~~~~~~~~~~~~~~~~~~~~~~~~~~~~~~~~~~~~~~~~~~~~~~

%% --------------------------------------------------------
\subsection*{Власна та домішкова провідність}
%% --------------------------------------------------------

У~чистому напівпровіднику електрони та дірки виникають парами. Якщо концентрація власних носіїв дорівнює $n_i$, то

% ========================================================
\begin{equation*}
	n=p=n_i.
\end{equation*}
% ========================================================


За підвищення температури частина електронів отримує достатню енергію для переходу через заборонену енергетичну зону,
і~концентрація носіїв швидко зростає.

У~наближенні невиродженого напівпровідника
% ~~~~~~~~~~~~~~~~~~~~~~~~~~~~~~~~~~~~~~~~~~~~~~~~~~~~~~~~
\begin{equation}\label{eq:IntrinsicConcentration}
n_i\propto T^{3/2}
\exp\left(-\frac{E_g}{2k_BT}\right),
\end{equation}
% ~~~~~~~~~~~~~~~~~~~~~~~~~~~~~~~~~~~~~~~~~~~~~~~~~~~~~~~~
де $E_g$~--- ширина забороненої зони.

Саме експоненційний множник є принципово важливим: навіть невелике підвищення температури може суттєво збільшити
концентрацію носіїв.

\begin{Attention}
Протилежні температурні залежності металів і~напівпровідників мають різне походження. У~металі температура головним чином
змінює \emphx{рухливість} уже наявних носіїв: $\tau$ зменшується, тому $\rho$ зростає. У~напівпровіднику температура
насамперед змінює \emphx{концентрацію} носіїв: $n$ та $p$ швидко зростають, тому $\rho$ спадає.
\end{Attention}

Електропровідність можна істотно змінити введенням невеликої кількості домішкових атомів. Домішки, які створюють додаткові
електрони, утворюють \emphx{напівпровідник $n$-типу}; домішки, що сприяють утворенню дірок,~--- \emphx{напівпровідник
$p$-типу}.

Для $n$-типу основними носіями є електрони, а~для $p$-типу~--- дірки; проте внесок неосновних носіїв завжди входить у~загальну формулу~\eqref{eq:SemiconductorConductivity}.
% ========================================================

%% --------------------------------------------------------
Типові донори в~кремнії~--- елементи V групи (фосфор, миш'як, сурма): п'ятий валентний електрон слабко зв'язаний з~атомом домішки (енергія іонізації $\approx0{,}045$~еВ), тому вже при кімнатній температурі практично всі донори іонізовані. Акцептори~--- елементи III групи (бор, алюміній, галій): вони захоплюють електрон з~валентної зони й~залишають дірку. Концентрації носіїв у~стані рівноваги пов'язані \emphx{законом діючих мас}
% ~~~~~~~~~~~~~~~~~~~~~~~~~~~~~~~~~~~~~~~~~~~~~~~~~~~~~~~~
\begin{equation}\label{eq:MassAction}
	np=n_i^2,
\end{equation}
% ~~~~~~~~~~~~~~~~~~~~~~~~~~~~~~~~~~~~~~~~~~~~~~~~~~~~~~~~
який справедливий і~при наявності домішок (домішки змінюють $n$ та $p$, але не їхній добуток). Для $n$-типу з~$N_d\gg n_i$ маємо $n\approx N_d$, $p=n_i^2/N_d$: дірки~--- \emphx{неосновні} носії, їх мізерно мало. Домішки змінюють провідність на багато порядків: один атом фосфору на кілька мільйонів атомів кремнію збільшує провідність у~$\sim10^5$ разів (задача~\ref{prob:DopedSilicon}).

Температурна залежність домішкового напівпровідника має три області (рис.~\ref{pic:DopedTemperature}). При дуже низьких $T$ донори <<виморожуються>>: електрони повертаються на домішкові рівні, і~$n$ експоненційно спадає. У~широкому інтервалі, що охоплює кімнатну температуру, усі донори іонізовані, $n\approx N_d=\const$, і~поведінка нагадує метал: рухливість спадає з~ростом $T$ (розсіяння на фононах, $\mu\propto T^{-3/2}$), тому $\rho$ \emphx{зростає}. Нарешті, при високих $T$ концентрація власних носіїв $n_i$ перевищує $N_d$, і~провідність знову експоненційно зростає. Саме тому кремнієві прилади працюють лише до $\sim150$--$200\,^\circ$C.

\begin{figure}[h!]\centering
	\localinput{DopedTemperature.tikz}
	% \figstub{DopedTemperature.tikz}{Графік $\ln n$ (або $\ln\lambda$) від $1/T$ для $n$-напівпровідника: виморожування, насичення (виснаження домішок), власна провідність}
	\caption{Три температурні області провідності домішкового напівпровідника\label{pic:DopedTemperature}}
\end{figure}

%% --------------------------------------------------------
\subsection*{Контакт двох напівпровідників}
%% --------------------------------------------------------

Якщо привести в~контакт напівпровідники $p$- та $n$-типів, виникає \emphx{$p$--$n$-перехід}. Електрони з~області $n$-типу
дифундують у~напрямку області $p$-типу, а~дірки~--- у~протилежному напрямку.

Поблизу межі контакту носії рекомбінують, і~утворюється збіднений шар, у~якому залишаються нерухомі заряджені домішки (рис.~\ref{pic:PNJunction}).
Ці заряди створюють внутрішнє електричне поле, яке перешкоджає подальшій дифузії основних носіїв.

\begin{figure}[h!]\centering
	\includegraphics[width=0.95\linewidth]{PNJunction}
	% \figstub{PNJunction.tikz}{Схема $p$--$n$-переходу: області $p$ і $n$, збіднений шар з нерухомими $\ominus$ та $\oplus$, напрямок внутрішнього поля, графіки $\rho(x)$, $E(x)$, $\phi(x)$ з висотою бар'єра $\phi_k$}
	\caption{Рівноважний $p$--$n$-перехід\label{pic:PNJunction}}
\end{figure}


У~результаті на межі виникає \emphx{контактна різниця потенціалів}. Зовнішня напруга може або посилювати внутрішнє поле,
або послаблювати його.

Якщо зовнішнє поле зменшує потенціальний бар'єр, через перехід проходить значний струм~--- це \emphx{пряме зміщення}.
При протилежному напрямку напруги бар'єр збільшується, і~струм основних носіїв різко зменшується~--- \emphx{зворотне зміщення}.

Саме ця асиметрія лежить в~основі роботи напівпровідникового \emphx{діода}.

Щоб зрозуміти, як працює перехід, розглянемо два механізми руху носіїв. Перший нам уже знайомий~--- \emphx{дрейф} у~полі, $\vect j=\lambda\Efield$. Другий~--- \emphx{дифузія}: носії прагнуть перейти з~області, де їх багато, в~область, де їх мало, навіть без поля. Рівновага (див. поле збідненого шару вище) настає, коли дрейфовий і~дифузійний струми взаємно скорочуються. Густина електронного струму
% ~~~~~~~~~~~~~~~~~~~~~~~~~~~~~~~~~~~~~~~~~~~~~~~~~~~~~~~~
\begin{equation}\label{eq:DriftDiffusion}
	\vect j_n = en\mu_n\Efield + eD_n\nabla n,
\end{equation}
% ~~~~~~~~~~~~~~~~~~~~~~~~~~~~~~~~~~~~~~~~~~~~~~~~~~~~~~~~
де $D_n$~--- коефіцієнт дифузії, пов'язаний з~рухливістю співвідношенням Ейнштейна $D_n=\mu_nk_BT/e$. Другий доданок відіграє ту саму роль, що й~стороння сила в~узагальненому законі Ома~\eqref{eq:GeneralizedOhm}: він створює струм не електричним полем, а~нерівномірністю концентрації.

Отже, $p$-$n$-перехід, як і~електролітична комірка (розд.~\ref{sec:CurrentInElectrolytes}),
--- це \emphx{неоднорідна ділянка кола}: окрім омічного опору, тут діє стороння сила
(дифузія носіїв), яка створює контактну різницю потенціалів $\phi_k$~\eqref{eq:ContactPotential}.
Саме тому до $p$--$n$-переходу незастосовний закон Ома у~формі $I=U/R$: він описується
формулою Шоклі~\eqref{eq:Shockley}, яка враховує цю сторонню силу.

У~рівновазі $\vect j_n=0$, звідки $\Efield=-\dfrac{k_BT}{e}\dfrac{\nabla n}{n}$. Інтегруючи поперек переходу, дістаємо \emphx{контактну різницю потенціалів}
% ~~~~~~~~~~~~~~~~~~~~~~~~~~~~~~~~~~~~~~~~~~~~~~~~~~~~~~~~
\begin{equation}\label{eq:ContactPotential}
	\phi_k=\frac{k_BT}{e}\ln\frac{n_n}{n_p}=\frac{k_BT}{e}\ln\frac{N_dN_a}{n_i^2},
\end{equation}
% ~~~~~~~~~~~~~~~~~~~~~~~~~~~~~~~~~~~~~~~~~~~~~~~~~~~~~~~~
де $n_n\approx N_d$ і~$n_p=n_i^2/N_a$~--- концентрації електронів у~$n$- та $p$-областях (закон діючих мас~\eqref{eq:MassAction}). Для кремнію $\phi_k\approx0{,}6$--$0{,}8$~В.

Прикладена напруга $U$ (додатна при прямому зміщенні) змінює висоту бар'єра до $e(\phi_k-U)$. Струм основних носіїв, що долають бар'єр, зростає як $e^{eU/k_BT}$, тоді як струм неосновних носіїв від $U$ практично не залежить. Результат описує формула Шоклі
% ~~~~~~~~~~~~~~~~~~~~~~~~~~~~~~~~~~~~~~~~~~~~~~~~~~~~~~~~
\begin{equation}\label{eq:Shockley}
	I=I_s\left(e^{eU/k_BT}-1\right),
\end{equation}
% ~~~~~~~~~~~~~~~~~~~~~~~~~~~~~~~~~~~~~~~~~~~~~~~~~~~~~~~~
де $I_s$~--- зворотний струм насичення (сильно, експоненційно залежить від температури). При прямому зміщенні $U\gg k_BT/e\approx26$~мВ струм зростає експоненційно, а~при зворотному $I\to-I_s$ (рис.~\ref{pic:PNIV}). Вольт-амперна характеристика $p$--$n$-переходу нелінійна, отже закон Ома для нього непридатний.

\begin{figure}[h!]\centering
	\localinput{PNIV.tikz}
	% \figstub{PNIV.tikz}{ВАХ діода за формулою Шоклі: різке зростання струму при прямому зміщенні, майже нуль ($-I_s$) при зворотному}
	\caption{Вольт-амперна характеристика $p$--$n$-переходу\label{pic:PNIV}}
\end{figure}

Застосування $p$--$n$-переходу: випрямні діоди; світлодіоди (рекомбінація електронів і~дірок при прямому зміщенні випромінює фотон); сонячні елементи (фотогенерація пар у~збідненому шарі); транзистори з~двома переходами, які лежать в~основі інтегральних схем.

\begin{Historical}
	Опір сульфіду срібла \ce{Ag2S} спадає при нагріванні~--- це помітив ще М.~Фарадей
	(1833). Зонну теорію напівпровідників розвинув А.~Вільсон (1931). У~Києві
	В.~Є.~Лашкарьов (1941) виявив $p$--$n$-перехід у~закису міді \ce{Cu2O}. Транзистор
	створили Дж.~Бардін і~В.~Браттейн (1947), теорію переходу дав В.~Шоклі (1949);
	усі троє~--- лауреати Нобелівської премії 1956~р.
\end{Historical}

%% --------------------------------------------------------
\section*{Підсумки}
\phantomsection
\addcontentsline{toc}{section}{Підсумки}
%% --------------------------------------------------------

Підсумуємо відмінності між середовищами (табл.~\ref{tab:CurrentCarriersMedia}): макроскопічна природа струму спільна, а~мікроскопічна реалізація різна.

\begin{table}[h!]\small\centering
\caption{Особливості електричного струму в~різних середовищах}
\label{tab:CurrentCarriersMedia}
\begin{tblr}{
colspec={X[l,m]X[l,m]X[l,m]},
row{odd} = {gray!10},
row{1} = {c, m, fg=white, bg=themecolorlight, font=\bfseries},
hlines,
vlines
}
Середовище      & Характерна особливість                                     & Закон Ома і залежність $\lambda(T)$                                                                   \\
Метали          & перенесення заряду без перенесення речовини                & виконується; $\lambda$ спадає з $T$ (менший $\tau$)                                                   \\
Електроліти     & струм супроводжується хімічними перетвореннями             & виконується; $\lambda$ зростає з $T$ (менша в'язкість)                                                \\
Гази            & концентрація носіїв залежить від іонізації та рекомбінації & лише в слабких полях; далі насичення та пробій                                                        \\
Вакуум          & струм можливий лише за наявності джерела носіїв            & не виконується (закон $U^{3/2}$, насичення); $j\propto T^2e^{-A_{\text{вих}}/k_BT}$                   \\
Напівпровідники & провідність сильно залежить від температури та домішок     & у однорідному зразку виконується, $\lambda$ експоненційно зростає з $T$; для $p$--$n$-переходу --- ні \\
\end{tblr}
\end{table}

Попри ці відмінності, для всіх середовищ залишається справедливим фундаментальний зв'язок~\eqref{eq:GeneralCurrentMedia}, тобто мікроскопічна основа поняття густини струму. Усі відмінності між провідниками зводяться до відповіді на три запитання:

\begin{enumerate}
\item \emphx{Які частинки є носіями заряду?}
\item \emphx{Скільки їх у~одиниці об'єму?}
\item \emphx{Як електричне поле змінює їхній рух?}
\end{enumerate}

У~цій главі закон Ома, що в~гл.~\ref{SteadyCurrent} постулювався як властивість середовища, набув мікроскопічного змісту: провідність $\lambda$ визначається тим, які носії є в~середовищі, скільки їх і~наскільки легко поле їх рухає~\eqref{eq:GeneralConductivity}.

За цією ознакою середовища природно діляться на дві групи. У~металах і~електролітах кількість носіїв задана самою речовиною, тому все визначає їхня рухливість, а~закон Ома виконується. У~газах, вакуумі й~напівпровідниках кількість носіїв залежить від зовнішніх умов~--- іонізації, емісії, температури, домішок,~--- і~зазвичай змінюється разом із полем, тому струм нелінійно залежить від напруги. Надпровідність є граничним випадком, де закон Ома втрачає зміст, бо $\rho=0$.

Усі середовища підкоряються загальним законам гл.~\ref{SteadyCurrent}: збереженню заряду, закону Джоуля-Ленца, узагальненому закону Ома зі сторонніми силами. Нові тут лише мікроскопічні причини, які визначають $\lambda$. Крім того, на прикладі металів видно межі класичної моделі: щоб правильно описати температурну залежність опору й~закон Відемана--Франца, довелося залучити принцип Паулі.
%% --------------------------------------------------------
\section*{Задачі}
\phantomsection
\addcontentsline{toc}{section}{Задачі}
%% --------------------------------------------------------
%=========================================================
\begin{problem}\label{prob:DrudeMeanFreeTime}
	\emphx{Час вільного пробігу за Друде.} Питома провідність міді $\lambda=5{,}9\cdot10^{7}$~См/м, концентрація електронів
	$n=8{,}5\cdot10^{28}$~м$^{-3}$. Оцініть середній час вільного пробігу $\tau$ за формулою Друде--Лоренца та \emphx{класичну} довжину
	вільного пробігу $\ell=v_T\tau$ при $T=300$~K ($v_T\approx1{,}2\cdot10^5$~м/с). Порівняйте одержане $\ell$ з~відстанню між іонами
	ґратки ($\sim0{,}3$~нм), а~також із реальною довжиною вільного пробігу в~міді при кімнатній температурі ($\sim40$~нм).
	% \emphx{Відповідь:} $\tau=\dfrac{m_e\lambda}{ne^2}\approx2{,}5\cdot10^{-14}$~с; $\ell\approx3$~нм --- це \emphx{класична} оцінка,
	% лише вдесятеро більша за міжіонну відстань. Експеримент же дає $\sim40$~нм, тобто на порядок більше: саме ця розбіжність і
	% спонукає до квантової теорії.
\end{problem}
%=========================================================
\begin{problem}\label{prob:CopperMobility}
	\emphx{Рухливість електронів у~міді.} Для міді ($\lambda=5{,}9\cdot10^{7}$~См/м, $n=8{,}5\cdot10^{28}$~м$^{-3}$) знайдіть рухливість $\mu$, напруженість поля, потрібну для густини струму $j=10^{6}$~А/м$^2$ (1~А через переріз 1~мм$^2$), і~дрейфову швидкість $u=\mu E$. Порівняйте з~результатом задачі~\ref{prob:DriftVelocityCopper} гл.~\ref{SteadyCurrent} і~виразіть $\mu$ через $\tau$ із задачі~\ref{prob:DrudeMeanFreeTime}.
	\emphx{Відповідь:} $\mu=\lambda/(en)\approx4{,}3\cdot10^{-3}$~м$^2$/(В$\cdot$с); $E=j/\lambda\approx1{,}7\cdot10^{-2}$~В/м; $u=\mu E\approx7\cdot10^{-5}$~м/с (збігається з~$u=I/(enS)$); $\mu=e\tau/m_e$ з~$\tau\approx2{,}5\cdot10^{-14}$~с дає те саме значення.
\end{problem}
%=========================================================
\begin{problem}\label{prob:TempCoefficient}
	\emphx{Температурний коефіцієнт опору.} При $20\,^\circ$C опір вольфрамового дроту $R_{20}=10$~Ом, а~$\alpha=5{,}0\cdot10^{-3}\,{}^\circ$C$^{-1}$.
	Яким буде опір при $t=1000\,^\circ$C, якщо вважати лінійну залежність~\eqref{eq:TempDependence2} справедливою? На скільки відсотків
	він зросте?
	\emphx{Відповідь:} $R=R_{20}[1+\alpha(t-20)]\approx59$~Ом; зростання на $\approx490\%$. (Зауважте: при $1000\,^\circ$C лінійна
	формула вже дає помітну похибку — реальний опір вольфраму трохи менший.)
\end{problem}

% (Стару версію нижче закоментовано: лінійна екстраполяція $\rho_{\text{фон}}\propto T$ до $1{,}8$~К фізично некоректна, бо при низьких $T$ $\rho_{\text{фон}}\propto T^5$.)
\begin{problem}\label{prob:MatthiessenRRR}
	\emphx{Правило Маттіссена.} Зразок міді має при $T=300$~К питомий опір $1{,}72\cdot10^{-8}$~Ом$\cdot$м, а~при $4{,}2$~К~--- $2\cdot10^{-10}$~Ом$\cdot$м (при $4{,}2$~К фононним внеском можна знехтувати, тож це $\rho_0$). Знайдіть фононний внесок при $300$~К, частку домішкового розсіяння та $\\rho(300\ \text{К})/\rho_0$. Яким буде $\rho(300\ \text{К})$ для чистішого зразка з~$\rho_0'=2\cdot10^{-11}$~Ом$\cdot$м?
	\emphx{Відповідь:} $\rho_{\text{фон}}=1{,}70\cdot10^{-8}$~Ом$\cdot$м; частка домішок $\approx1{,}2\%$; $\text{RRR}\approx86$; для чистішого зразка $\rho(300\ \text{К})\approx1{,}702\cdot10^{-8}$~Ом$\cdot$м~--- при кімнатній температурі чистота майже не впливає, на відміну від низьких температур.
\end{problem}

%\begin{problem}\label{prob:Matthiessen}
%	\emphx{Правило Маттіссена.} Зразок міді має залишковий опір $\rho_0=1{,}0\cdot10^{-10}$~Ом$\cdot$м і при $T=300$~K повний питомий опір
%	$\rho=1{,}7\cdot10^{-8}$~Ом$\cdot$м. Знайдіть фононний внесок $\rho_{\text{фон}}$ і оцініть, за якої температури він зрівняється із
%	залишковим (вважайте $\rho_{\text{фон}}\propto T$).
%	\emphx{Відповідь:} $\rho_{\text{фон}}\approx1{,}69\cdot10^{-8}$~Ом$\cdot$м; рівність при $T\approx1{,}8$~К.
%\end{problem}

\begin{problem}\label{prob:WiedemannFranz}
	\emphx{Число Лоренца.} Для міді при $T=300$~K питома електропровідність $\lambda=5{,}9\cdot10^{7}$~См/м, а~теплопровідність
	$\kappa=4{,}0\cdot10^{2}$~Вт/(м$\cdot$К). Знайдіть відношення $\kappa/\lambda$ і~порівняйте його з~добутком $LT$ для класичного~\eqref{eq:LorenzClassical} та квантового~\eqref{eq:LorenzQuantum} значень числа Лоренца. Яка з~двох теорій узгоджується з~дослідом?
	Візьміть $k_B=1{,}38\cdot10^{-23}$~Дж/К, $e=1{,}60\cdot10^{-19}$~Кл.
	% \emphx{Відповідь:} $\kappa/\lambda\approx6{,}8\cdot10^{-6}$~Вт$\cdot$Ом/К; $L_{\text{кл}}T\approx3{,}3\cdot10^{-6}$, $LT\approx7{,}3\cdot10^{-6}$~Вт$\cdot$Ом/К ---
	% з квантовим значенням збіг у межах $\sim10\%$, класичне вдвічі занижене.
\end{problem}

\begin{problem}\label{prob:SuperconductingWire}
	\emphx{Надпровідний дріт.} Волокна \ce{NbTi} у~дроті мають сумарний переріз $S=0{,}5$~мм$^2$ і~критичну густину струму $j_c=2\cdot10^{9}$~А/м$^2$. (а) Який максимальний струм може нести дріт без переходу в~нормальний стан? (б) Яка потужність виділялась би на $1$~м довжини в~мідному дроті того самого перерізу при цьому струмі ($\rho=1{,}68\cdot10^{-8}$~Ом$\cdot$м)? (в) Які матеріали з~розділу про надпровідність можна охолоджувати рідким азотом ($77$~К)?
	\emphx{Відповідь:} (а) $I_c=j_cS=1000$~А; (б) $P/\ell=I^2\rho/S\approx34$~Вт/м (закон Джоуля--Ленца~\eqref{eq:JouleLenzIntegral}), у~надпровіднику~--- нуль; (в) ті, у~яких $T_c>77$~К: \ce{YBa2Cu3O7}, Bi-купрат, Hg-купрат (\ce{H3S}~--- лише під тиском $\sim150$~ГПа).
\end{problem}

\begin{problem}\label{prob:ElectrolysisCopper}
	\emphx{Мідне покриття.} Крізь розчин мідного купоросу \ce{CuSO4} з~мідними електродами пропускають струм $I=2$~А протягом
	$t=30$~хв. Яка маса міді виділиться на катоді? Молярна маса міді $\mu=63{,}5$~г/моль, валентність $z=2$, стала Фарадея
	$F=9{,}65\cdot10^{4}$~Кл/моль.
	\emphx{Відповідь:} $m=\dfrac{\mu}{zF}It\approx1{,}18$~г.
\end{problem}

\begin{problem}\label{prob:ElectrolyteConductivity}
	\emphx{Розчин кухонної солі.} У~розчині \ce{NaCl} концентрація йонів кожного знака $n_\pm=6{,}0\cdot10^{25}$~м$^{-3}$ (дисоціацію
	вважайте повною), рухливості $\mu_+=5{,}2\cdot10^{-8}$~м$^2$/(В$\cdot$с) для \ce{Na+} і~$\mu_-=7{,}9\cdot10^{-8}$~м$^2$/(В$\cdot$с)
	для \ce{Cl-}. Знайдіть питому провідність розчину та частку струму, яку переносять катіони.
	\emphx{Відповідь:} $\lambda=en(\mu_++\mu_-)\approx1{,}3$~См/м; частка катіонів $\mu_+/(\mu_++\mu_-)\approx0{,}40$.
\end{problem}

\begin{problem}\label{prob:ElectrolysisPolarization}
	\emphx{Електроліз води.} На електроди ванни подано $U=2{,}0$~В, ЕРС поляризації $\emf_{\text{пол}}=1{,}5$~В, опір розчину $R=0{,}05$~Ом. Знайдіть струм, масу водню за $t=1$~год ($\mu=2{,}016$~г/моль, $z=2$ для молекули \ce{H2}), потужність тепловиділення в~розчині та частку потужності джерела, що йде на хімічну реакцію.
	\emphx{Відповідь:} $I=(U-\emf_{\text{пол}})/R=10$~А (формула~\eqref{eq:ElectrolysisOhm}); $m=\dfrac{\mu}{zF}It\approx0{,}38$~г; $P_{\text{Дж}}=I^2R=5$~Вт; частка $\emf_{\text{пол}}/U=75\%$.
\end{problem}


\begin{problem}\label{prob:SaturationCurrent}
	\emphx{Струм насичення.} Іонізатор створює в~газі між плоскими електродами площею $S=100$~см$^2$, розташованими на відстані
	$d=2$~см один від одного, $N=10^{12}$ пар йонів в~одиниці об'єму за одиницю часу (м$^{-3}\cdot$с$^{-1}$). Знайдіть струм насичення.
	\emphx{Відповідь:} $I_{\text{нас}}=eNSd\approx3{,}2\cdot10^{-11}$~А.
\end{problem}

\begin{problem}\label{prob:IonizedAirConductivity}
	\emphx{Іонізоване повітря.} Іонізатор створює в~повітрі $N=10^{13}$~м$^{-3}\cdot$с$^{-1}$ пар йонів; коефіцієнт рекомбінації $\beta=1{,}6\cdot10^{-12}$~м$^3$/с, рухливості йонів $\mu_+=1{,}4\cdot10^{-4}$ та $\mu_-=1{,}9\cdot10^{-4}$~м$^2$/(В$\cdot$с). Знайдіть рівноважну концентрацію йонів і~питому провідність. Оцініть напругу на електродах, розташованих на відстані $d=2$~см, до якої виконується закон Ома.
	\emphx{Відповідь:} $n_0=\sqrt{N/\beta}=2{,}5\cdot10^{12}$~м$^{-3}$; $\lambda=en_0(\mu_++\mu_-)\approx1{,}3\cdot10^{-10}$~См/м; $j_{\text{нас}}=eNd\approx3{,}2\cdot10^{-8}$~А/м$^2$, $E\sim j_{\text{нас}}/\lambda\approx2{,}4\cdot10^{2}$~В/м, $U\sim Ed\approx5$~В.
\end{problem}



\begin{problem}\label{prob:ImpactIonization}
	\emphx{Ударна іонізація.} Енергія іонізації атома аргону $W_i=15{,}8$~еВ. Яка мінімальна напруженість однорідного поля потрібна,
	щоб електрон, який пролітає без зіткнень відстань $\ell=0{,}1$~мм, набув на ній енергії, достатньої для іонізації? Втратами енергії
	на цьому шляху знехтуйте.
	\emphx{Відповідь:} $E=\dfrac{W_i}{e\ell}\approx1{,}6\cdot10^{5}$~В/м.
\end{problem}

\begin{problem}\label{prob:CoronaOnset}
	\emphx{Початок корони.} Електрична міцність повітря $E_{\text{пр}}=3\cdot10^{6}$~В/м. Знайдіть потенціал ізольованої провідної кулі, при якому біля її поверхні починається коронний розряд, для $r=1$~см та $r=0{,}1$~мм (поле біля поверхні $E=\phi/r$). Чому корона виникає біля вістер?
	\emphx{Відповідь:} $\phi=E_{\text{пр}}r=30$~кВ і~$300$~В відповідно; на вістрі (малий радіус кривизни) поле при тому самому потенціалі найбільше.
\end{problem}



\begin{problem}\label{prob:RichardsonRatio}
	\emphx{Нагрів катода.} Робота виходу вольфраму $A_{\text{вих}}=4{,}5$~еВ. У~скільки разів зросте густина термоелектронного струму при нагріванні
	катода від $T_1=2000$~К до $T_2=2500$~К? Скористайтеся законом Річардсона--Дешмана~\eqref{eq:RichardsonDushman}; стала Больцмана
	$k_B=8{,}62\cdot10^{-5}$~еВ/К.
	\emphx{Відповідь:} $\dfrac{j_2}{j_1}=\left(\dfrac{T_2}{T_1}\right)^{2}\exp\!\left[\dfrac{A_{\text{вих}}}{k_B}\left(\dfrac1{T_1}-\dfrac1{T_2}\right)\right]\approx2{,}9\cdot10^{2}$.
\end{problem}

% Зауваження: ця задача --- чиста кінематика і сама по собі слабко пов'язана з текстом; залишена як розминка перед задачею ChildLangmuir.

\begin{problem}\label{prob:DiodeElectronSpeed}
	\emphx{Електрон у~діоді.} Електрон вилітає з~катода вакуумного діода з~нехтовно малою швидкістю й~прискорюється напругою
	$U=200$~В. Якою буде його швидкість біля анода і~за який час він пролетить проміжок $d=5$~мм, якщо поле в~ньому вважати однорідним?
	\emphx{Відповідь:} $v=\sqrt{2eU/m_e}\approx8{,}4\cdot10^{6}$~м/с; $t=2d/v\approx1{,}2\cdot10^{-9}$~с.
\end{problem}

\begin{problem}\label{prob:ChildLangmuir}
	\emphx{Струм, обмежений просторовим зарядом.} Плоский вакуумний діод має відстань між електродами $d=5$~мм, напругу $U=200$~В, площу електродів $S=1$~см$^2$; катод емітує необмежену кількість електронів. Знайдіть густину і~силу струму за законом $3/2$~\eqref{eq:ChildLangmuir}. Як зміняться $j$ при подвоєнні $U$ та при подвоєнні $d$? Яка потужність виділяється на аноді?
	\emphx{Відповідь:} $j=\dfrac{4\epsilon_0}{9}\sqrt{\dfrac{2e}{m_e}}\dfrac{U^{3/2}}{d^2}\approx2{,}6\cdot10^{2}$~А/м$^2$; $I=jS\approx26$~мА; $j$ зросте у~$2^{3/2}\approx2{,}8$ раза при $2U$ і~зменшиться вчетверо при $2d$; $P=IU\approx5$~Вт.
\end{problem}


\begin{problem}\label{prob:IntrinsicSilicon}
	\emphx{Власний кремній.} При $T=300$~К концентрація власних носіїв у~кремнії $n_i=1{,}5\cdot10^{16}$~м$^{-3}$, а~рухливості
	електронів і~дірок $\mu_n=0{,}135$~м$^2$/(В$\cdot$с) та $\mu_p=0{,}048$~м$^2$/(В$\cdot$с). Знайдіть питому провідність і~питомий
	опір власного кремнію.
	\emphx{Відповідь:} $\lambda=en_i(\mu_n+\mu_p)\approx4{,}4\cdot10^{-4}$~См/м; $\rho\approx2{,}3\cdot10^{3}$~Ом$\cdot$м.
\end{problem}

\begin{problem}\label{prob:IntrinsicTemperature}
	\emphx{Нагрів напівпровідника.} У~скільки разів зросте концентрація власних носіїв у~кремнії ($E_g=1{,}12$~еВ) при нагріванні
	від $T_1=300$~К до $T_2=330$~К? Порівняйте із зміною питомого опору міді ($\alpha=4{,}3\cdot10^{-3}\,{}^\circ$C$^{-1}$) за тієї
	самої зміни температури. Зміною рухливостей знехтуйте.
	\emphx{Відповідь:} $\dfrac{n_i(T_2)}{n_i(T_1)}=\left(\dfrac{T_2}{T_1}\right)^{3/2}\exp\!\left[\dfrac{E_g}{2k_B}\left(\dfrac1{T_1}-\dfrac1{T_2}\right)\right]\approx8$, тобто
	опір кремнію \emphx{спадає} майже у~$8$ разів, тоді як опір міді \emphx{зростає} лише на $\alpha\Delta t\approx13\%$.
\end{problem}

\begin{problem}\label{prob:DopedSilicon}
	\emphx{Легований кремній.} У~кремній ввели донорну домішку з~концентрацією $N_d=10^{22}$~м$^{-3}$. При $T=300$~К вважайте всі
	донори іонізованими, тож $n\approx N_d$. Знайдіть питому провідність, концентрацію дірок (закон діючих мас $np=n_i^2$) і~в~скільки
	разів провідність перевищує провідність власного кремнію (значення $n_i$, $\mu_n$, $\mu_p$ візьміть з
	задачі~\ref{prob:IntrinsicSilicon}).
	\emphx{Відповідь:} $\lambda\approx en\mu_n\approx2{,}2\cdot10^{2}$~См/м; $p=n_i^2/n\approx2\cdot10^{10}$~м$^{-3}$ (дірки~--- неосновні носії);
	провідність зросла приблизно в~$5\cdot10^{5}$ разів.
\end{problem}


\begin{problem}\label{prob:PNJunction}
	\emphx{$p$--$n$-перехід у~кремнії.} Для кремнію ($n_i=1{,}5\cdot10^{16}$~м$^{-3}$) при $T=300$~К ($k_BT/e\approx25{,}9$~мВ) з~$N_a=10^{23}$~м$^{-3}$ та $N_d=10^{22}$~м$^{-3}$ знайдіть контактну різницю потенціалів~\eqref{eq:ContactPotential}. У~скільки разів прямий струм при $U=+0{,}5$~В перевищує модуль зворотного при $U=-0{,}5$~В (формула Шоклі~\eqref{eq:Shockley})?
	\emphx{Відповідь:} $\phi_k=\dfrac{k_BT}{e}\ln\dfrac{N_aN_d}{n_i^2}\approx0{,}75$~В; $\dfrac{I(+0{,}5)}{|I(-0{,}5)|}\approx e^{19{,}3}\approx2{,}5\cdot10^{8}$.
\end{problem}

\begin{Questions}
\begin{quizlist}
	\item Запишіть загальний вираз густини струму через сорти носіїв, $\vect j=\sum_k n_kq_k\vect u_k$, і~виведіть питому провідність $\lambda=\sum_k n_k|q_k|\mu_k$. Що називають рухливістю носіїв?
	\item Опишіть модель провідності металів Друде: як з~рівняння руху електрона між зіткненнями отримують питому провідність $\lambda = ne^2\tau/(2m_e)$? Чому множник $1/2$ у~цій формулі~--- наслідок спрощеного припущення?
	\item Як залежить опір металів від температури і~чим ця залежність відрізняється від поведінки напівпровідників? Що таке надпровідність?
	\item Чому класична теорія провідності металів (навіть у~варіанті Друде--Лоренца) не може пояснити лінійну залежність $\rho\propto T$, залишковий опір $\rho_0$ і~великі довжини вільного пробігу при низьких температурах? Сформулюйте правило Маттіссена.
	\item Що таке надпровідність і~критична температура? Чому струм у~надпровідному кільці може текти без джерела ЕРС і~як це узгоджується із законом Джоуля-Ленца? Поясніть роль куперівських пар у~теорії БКШ і~значення відкриття високотемпературної надпровідності 1986~р.
	\item Сформулюйте закон Відемана--Франца. Які значення числа Лоренца дають класична й~квантова теорії? Чому збіг Друде з~експериментом~--- радше щасливий випадок (компенсація помилок), ніж підтвердження теорії?
	\item Чому мідь, срібло й~золото~--- чудові провідники, але не надпровідники? Як це пов'язано з~механізмом утворення куперівських пар?
	\item Які носії струму в~електроліті? Чому внески катіонів і~аніонів у~густину струму додаються, хоча рухаються вони в~протилежних напрямках?
	\item Сформулюйте закони Фарадея для електролізу й~виведіть формулу $m=\mu It/(zF)$. Що таке стала Фарадея?
	\item Що таке іонізація й~рекомбінація? Чому концентрація носіїв у~газі, на відміну від металу, не є сталою?
	\item Опишіть хід залежності струму від напруги при несамостійному газовому розряді. Що таке струм насичення і~від чого він залежить?
	\item Що таке ударна іонізація й~електронна лавина? Чим самостійний розряд відрізняється від несамостійного? Назвіть основні види самостійного розряду.
	\item Які види електронної емісії ви знаєте? Що таке робота виходу і~як густина термоелектронного струму залежить від температури?
	\item Чому вакуумний діод має властивість односторонньої провідності? Чи є в~цій ситуації <<струм вакууму>> властивістю самого простору?
	\item Що таке дірка? Чому дірку не можна вважати самостійною елементарною частинкою?
	\item Чому питомий опір напівпровідника зазвичай спадає з~ростом температури, а~металу~--- зростає? Порівняйте роль концентрації та рухливості носіїв в~обох випадках.
	\item Чим відрізняються напівпровідники $n$- і~$p$-типу? Як домішки змінюють провідність і~які носії при цьому є основними?
	\item Як виникає $p$--$n$-перехід і~чому він проводить струм лише в~одному напрямку? Що таке пряме й~зворотне зміщення?
	\item На які три запитання зводяться всі відмінності між провідними середовищами? Дайте відповіді для металу, електроліту, газу, вакууму й~напівпровідника.
	\item Що показує сфера Фермі в~електричному полі? Чим змінюється зміст часу $\tau$ у~квантовій теорії провідності й~чому $\rho_{\text{фон}}\propto T$ при $T\gtrsim\Theta_D$ та $\propto T^5$ при низьких $T$?
	\item Запишіть закон Ома для електролітичної комірки. Що таке ЕРС поляризації й~чим електроліз відрізняється від роботи гальванічного елемента?
	\item Запишіть баланс іонізації й~рекомбінації, виведіть струм насичення $j_{\text{нас}}=eNd$. Сформулюйте умову Таунсенда й~закон Пашена.
	\item Виведіть закон трьох других для вакуумного діода. Чому тут закон Ома не виконується?
	\item Сформулюйте закон діючих мас. Опишіть три температурні області провідності домішкового напівпровідника.
	\item Як дифузія виконує роль сторонніх сил у~$p$-$n$-переході? Звідки береться контактна різниця потенціалів і~формула Шоклі?
\end{quizlist}
\end{Questions}